\documentclass[10pt,a4paper]{amsart}
\usepackage{amsmath,amssymb,amsthm,graphicx}
\usepackage{hyperref}
\usepackage{amsmath}
\usepackage{amssymb}
\usepackage{amsthm}
\usepackage{amscd}
\usepackage{graphicx}
\usepackage{varioref}
\numberwithin{equation}{section}

\newtheorem{Theorem}{Theorem}[section]
\newtheorem*{Theorem*}{Theorem}
\newtheorem{Corollary}[Theorem]{Corollary}
\newtheorem{Lemma}[Theorem]{Lemma}
\newtheorem{Proposition}[Theorem]{Proposition}
 { \theoremstyle{definition}
\newtheorem{Definition}[Theorem]{Definition}
\newtheorem{Note}[Theorem]{Note}
\newtheorem{Example}[Theorem]{Example}
\newtheorem{Remark}[Theorem]{Remark} }
\allowdisplaybreaks
\pdfstringdefDisableCommands{
\def\boldsymbol#1{#1}
\def\operatorname#1{#1}
\def\overline#1{#1}
}
\title{Compactifiable deformation families of complete quasi-projective varieties and symmetric projective K3 surfaces}
\author{Fan Xu}
\date{doi:10.3842/SIGMA.2025.062; CC BY 4.0}
\begin{document}
\maketitle
\noindent\textit{Allowed standard prerequisites (editorial).}\begin{itemize}
\item Research prerequisites named in the original article. Full original bibliography retained; all source-local definitions and supporting arguments retained.
\end{itemize}\medskip
\section{Introduction}

This paper is motivated by the gluing construction of K3 surfaces through two open complements of the closures of tubular neighborhoods of elliptic curves embedded in \smash{$\operatorname{CP}^2 \#9\overline{\operatorname{CP}}^2$} in \cite{Takayuki-Takato1}. Moreover, a gluing construction of projective K3 surfaces was presented in \cite{Takayuki-Takato2}.

For a rational elliptic surface $ X $ defined as the blow-up of a projective plane at nine base points for a pencil of cubics, Hans-Joachim Hein presented some complete Calabi--Yau metrics on $ X \backslash D $ with $D$ being a fiber on $X$ in \cite{Hans}. However, the blow up of a projective plane at arbitrary nine points may not be a rational elliptic surface.

Let $ \langle 1,\tau \rangle$ denote a lattice for $\tau\in Y:=\{\tau \mid \tau\in \mathbb{C},\, \operatorname{\operatorname{Im}}\tau>0 \}$. Then $C_0(\tau)=\mathbb{C}/ \langle 1,\tau \rangle$ is an elliptic curve for $\tau\in Y$. Let $C_0=\mathbb{C}/ \langle 1,\tau_0 \rangle$ with $\tau_0 \in Y$. The smooth elliptic curve $C_0 (\tau)$ can be holomorphically embedded in $\operatorname{CP}^2$ for $\tau \in Y$. The embedding map is induced by the Weierstrass $\wp$-function
\[\wp(z) =\sum _{\lambda \in \langle 1,\tau \rangle \backslash \{0\}} \left(\frac{1}{(z-\lambda )^2}-\frac{1}{\lambda ^2}\right)+\frac{1}{z^2}\]
and the Eisenstein series
\[G_{2k} ( \langle 1,\tau \rangle)=\sum _{\lambda \in \langle 1,\tau \rangle \backslash \{0\}} \lambda ^{-2k}.\]
Here, for $\tau\in Y$, the holomorphic embedding map $f_\tau\colon C_0 (\tau) \rightarrow\operatorname{CP}^2$ is defined by
\[f_\tau([z])=[\wp(z):\wp'(z):1]=[z_1:z_2:z_3]\]
for $[z]\in C_0 (\tau)$ with $[z]\neq[0]$ and $f_\tau([0])=[0:1:0]$. Furthermore,
\[f_\tau(C_0(\tau))=\bigl\{ [z_1:z_2:z_3]|-4 z_1^3+60 G_4 ( \langle 1,\tau \rangle)z_3^2 z_1+140 G_6 ( \langle 1,\tau \rangle)z_3^3+z_2^2 z_3=0\bigr\}\subset\operatorname{CP}^2\]
for $\tau \in Y$.

So $f_\tau(C_0(\tau))$ is a submanifold of $S_0=\operatorname{CP}^2$. Let \smash{$S(\tau)\cong\operatorname{CP}^2 \#9\overline{\operatorname{CP}}^2$} be the blow up of $S_0$ at nine points in the set $Z:=\{p_1,p_2,\dots ,p_9\}$ on $f_\tau(C_0(\tau))$. Now let $C(\tau)$ denote the strict transform of $f_\tau(C_0(\tau))$ in $S(\tau)$. Then $S(\tau)\setminus C(\tau)$ is a quasi-projective variety for $\tau\in Y$. In particular, let $S=S(\tau_0)$ and $C=C(\tau_0)$. Here, the normal bundle of $C(\tau)$ in $S(\tau)$ is assumed to satisfy the Diophantine condition corresponding to a Diophantine number pair proposed in~\cite{Takayuki-Takato2}. Then a complex analytic family of \smash{$\operatorname{CP}^2 \#9\overline{\operatorname{CP}}^2$} over a 9-dimensional complex manifold is constructed exactly. The simple description of this deformation family can be found in \cite{Takayuki-Takato1}.

Elizabeth Gasparim and Francisco Rubilar presented new definitions for deformation family in \cite{Elizabeth-Francisco}, especially for the deformation family of open manifolds. Taking use of the new definitions, the main theorem is as follows.


\begin{Theorem} \label{Theorem1.1}
There is a smooth compactifiable deformation of $S\backslash C$ over a $9$-dimensional complex manifold $T$. In addition, the deformation is differentially trivial along $T$.
\end{Theorem}


As it was proved in \cite{Anna}, $\forall \tau_1, \tau_2 \in\! Y$, if there was an algebraic isomorphism between $S(\tau_1)\backslash C(\tau_1)$ and $S(\tau_2)\backslash C(\tau_2)$, then there was an induced birational morphism between $S(\tau_1)$ and $S(\tau_2)$ which made the square-zero elliptic curves $C(\tau_1)$ and $C(\tau_2)$ isomorphic to each other. Therefore, ${\forall \tau_1 \neq \tau_2 \in Y}$, if $C(\tau_1)$ is not isomorphic to $C(\tau_2)$, then $S(\tau_1)\backslash C(\tau_1)$ is not algebraically isomorphic to $S(\tau_2)\backslash C(\tau_2)$. So through suitable choice of $T$, every two different fibers of the compactifiable deformation of $S\backslash C$ in Theorem~\ref{Theorem1.1} are not algebraically isomorphic to each other.

A K3 surface with an involution was constructed as an example in \cite[Section 7.1.2]{Takayuki-Takato1} by Takayuki Koike and Takato Uehara. Then through fixing suitable ample line bundles on $S(\tau)$, a gluing construction of symmetric projective K3 surfaces will be introduced here.

Firstly, there are ample line bundles properly defined on each fiber of the deformation family of \smash{$\operatorname{CP}^2 \#9\overline{\operatorname{CP}}^2$} \cite{Tomasz-Halszka}.

Secondly, in \cite{Takayuki-Takato2}, there was a simple description for the construction of some complete K\"ahler metrics on $\text{S}(\tau)\backslash \text{C}(\tau)$ corresponding to an ample line bundle on $S(\tau)$. In this paper, through modifying the method proposed in \cite{Takayuki-Takato2}, the analogous complete K\"ahler metrics corresponding to the ample line bundles selected are constructed on each fiber of the smooth compactifiable deformation family of $S \backslash C$.

Then the smooth compactifiable deformation family is proved to be a deformation family of complete quasi-projective varieties.

In particular, Yoshio Fujimoto proposed a condition for the nine blowing up points such that~$S(\tau)$ was a rational elliptic surface in \cite{Yoshio}. Taking use of the same construction method as Section~\ref{sec3} of this paper, we may construct another compactifiable deformation family and each fiber is an open Calabi--Yau manifold. The Calabi--Yau metrics on the fibers were constructed in~\cite{Hans} by Hans-Joachim Hein. However, this situation will not be further discussed in this paper.

Thirdly, a symmetric projective K3 surface can be constructed through gluing the open completement of closure of the tubular neighborhood of $C(\tau)$ in $S(\tau)$ and an analogous manifold.

Then a deformation family of symmetric projective K3 surfaces is constructed as follows.


\begin{Corollary} \label{Corollary1.2}
A deformation family of symmetric projective K3 surfaces over a $10$-dimensional complex manifold along with a symmetric K\"ahler metric on each fiber can be constructed through a smooth compactifiable deformation of $S \backslash C$ and an analogous deformation family.
\end{Corollary}


Moreover, the complex analytic family of symmetric projective K3 surfaces has an injective Kodaira--Spencer map by \cite[Theorem~1.1]{Takayuki-Takato1}.

The following is to introduce the main contents in the following sections of this paper.

In Section~\ref{sec2}, the main goal is to present a gluing construction of symmetric projective K3 surfaces and the toroidal groups corresponding to Diophantine number pairs.

Firstly, the Diophantine condition is introduced for the construction of a tubular neighborhood of an elliptic curve embedded in \smash{$\operatorname{CP}^2 \#9 \overline{\operatorname{CP}}^2$}.

Secondly, fixing a suitable ample line bundle on $S(\tau)$, the open completement of the closure of the tubular neighborhood of $C(\tau)$ in $S(\tau)$ and an analogous open manifold can be glued to a~symmetric projective K3 surface through a suitable map.

Finally, the specific defined toroidal groups related to Diophantine number pairs are just properly defined in a special class of toroidal groups described in \cite[Theorem]{Christian}.

In Section~\ref{sec3}, the main goal is to prove Theorem~\ref{Theorem1.1} and construct complete K\"ahler metrics on each quasi-projective variety fiber of the deformation family through fixing an ample line bundle on $S(\tau)$.

Moreover, through the process of constructing the complete K\"ahler metrics on $S(\tau)\backslash C(\tau)$, the condition for the gluing construction of symmetric projective K3 surfaces can be confirmed.

In Section~\ref{sec4}, the main goal is to prove Corollary~\ref{Corollary1.2}.

\section{Symmetric projective K3 surfaces and toroidal groups}\label{sec2}

\subsection{A gluing construction of symmetric projective K3 surfaces}

Firstly, it is to give the details for a gluing construction of symmetric projective K3 surfaces.

Julius Ross and David Witt Nystr\"om stated that there was a canonical smooth tubular neighborhood for a compact complex submanifold of a K\"ahler manifold which in general would not be holomorphic in \cite{Julius-David}.

V.I.~Arnold proposed a theorem that in most situations an elliptic curve holomorphically embedded in a complex surface with self-intersection index 0 should have a small neighborhood biholomorphically equivalent to a neighborhood of zero section of the normal bundle in \cite{Vladimir}. Ueda also mentioned a corresponding theorem in \cite{Tetsuo}.

Let $\mathbb{N} \subset \mathbb{Z} $ denote the set consisting of positive integers.

Another fact is that any topologically trivial holomorphic line bundle on an elliptic curve is flat \cite{Tetsuo}. The following is a smooth version definition for the monodromy of a complex line bundle over a Riemann surface.

\begin{Definition} \label{Definition2.1}
Let $\nabla$ be a flat connection on an arbitrary complex line bundle $L$ representing an isomorphism class of complex line bundles over a Riemann surface $M$. Define $g_0={\rm e}^{\sqrt{-1} t}\in S^1$ for $t\in [0,2 \pi ]$ and let $\gamma \colon S^1\to M$ be a smooth closed curve. In addition, let $\psi$ denote a parallel section of $\widehat{\gamma}^*L$ with $\widehat{\gamma}=\gamma \circ g_0 $. Then $\widehat{f}(\gamma)\in S^1$ with $\psi_{2\pi} =\widehat{f}(\gamma)\psi_0$ is defined to be the monodromy of $L$ along $\gamma$.
\end{Definition}

The Diophantine condition for a pair of real numbers from \cite{Takayuki-Takato2} is defined as follows.

\begin{Definition}[{\cite[Definition~2.1]{Takayuki-Takato2}}] \label{Definition2.2}
A pair of numbers $(p,q)\in \mathbb{R}^2$ is said to satisfy the Diophantine condition if there exist $\vartheta >0$ and $A>0$ such that
\[
\underset{\mu_0 ,\nu_0 \in \mathbb{Z}}{\min } \big| n\bigl(p+q\sqrt{-1} \bigr)-\bigl(\mu_0 +\nu_0\sqrt{-1} \bigr)\big| \geq A\cdot n^{-\vartheta }
\]
for any $n \in \mathbb{N}$.
\end{Definition}

With the definitions above, the Diophantine condition for a flat holomorphic line bundle on an elliptic curve embedded in a complex surface for the construction of a holomorphic tubular neighborhood of the elliptic curve is defined as follows.

\begin{Definition}[{\cite[Definition~2.2]{Takayuki-Takato2}}] \label{Definition2.3}
For $\tau\in Y$, let $\widetilde{\alpha}$ and $\widetilde{\beta}$ be smooth loops on the elliptic curve $C_0(\tau)=\mathbb{C}/ \langle 1,\tau \rangle$ corresponding to the line segments $[0,1]$ and $[0,\tau]$, respectively. For any topologically trivial line bundle $L\in \operatorname{Pic}^0(C_0(\tau))$, let \smash{${\rm e}^{\widetilde{p}\cdot2\pi\sqrt{-1}}$} and \smash{${\rm e}^{\widetilde{q}\cdot2\pi\sqrt{-1}}$} denote the monodromies of $L$ along the loops $\widetilde{\alpha}$ and \smash{$\widetilde{\beta}$} respectively with $(\widetilde{p},\widetilde{q})\in \mathbb{R}^2$. Then $L$ is said to satisfy the Diophantion condition if $(\widetilde{p},\widetilde{q})$ satisfies the Diophantine condition.
\end{Definition}

According to \cite[Theorem~1.6]{Takayuki-Takato1}, if $\widetilde{L}=N_{C(\tau)/S(\tau)}\in \operatorname{Pic}^0(C(\tau))$ satisfies the Diophantine condition in Definition~\ref{Definition2.3}, \smash{$-\log d\bigl(\mathbb{I}_{C(\tau)},\widetilde{L}^n\bigr)=O(\log n)$} as $n\rightarrow\infty$ with $d$ being an invariant distance on $\operatorname{Pic}^0(C(\tau))$ and $\mathbb{I}_{C(\tau)}$ being the holomorphically trivial line bundle on $C(\tau)$. So from~\cite[Theorem 3]{Tetsuo} which is a generalized result originating from \cite{Vladimir} by V.I.~Arnold, the following theorem is true.

\begin{Theorem}[\cite{Takayuki-Takato2}] \label{Theorem2.4}
If the normal bundle \emph{$N_{C(\tau)/S(\tau)}\in \operatorname{Pic}^0(C(\tau))$} satisfies the Diophantine condition, then $C(\tau)$ admits a holomorphic tubular neighborhood in \smash{\emph{$S(\tau)\cong\operatorname{CP}^2\text{$\#$9}\overline{\operatorname{CP}}^2$}} biholomorphic to a neighborhood of the zero section in \smash{$N_{C(\tau)/S(\tau)}$}.
\end{Theorem}

Now assume that \smash{$N_{C(\tau)/S(\tau)}$} satisfies the Diophantine condition. Let $\alpha$ and $\beta$ be smooth loops on $C(\tau)$ corresponding to the line segments $[0,1]$ and $[0,\tau]$ on $C_0(\tau)$. Let \smash{${\rm e}^{p\cdot2\pi\sqrt{-1}}$} and~\smash{${\rm e}^{q\cdot2\pi\sqrt{-1}}$} be the monodromies of $N_{C(\tau)/S(\tau)}$ along the loops $\alpha$ and $\beta$ respectively with $(p,q)\in \mathbb{R}^2$. So $(p,q)$ should satisfy the Diophantine condition. Taking use of Theorem~\ref{Theorem2.4}, the following statement is true.

There is a tubular neighborhood $W$ of $C(\tau)$ biholomorphic to a neighborhood of the zero section in $N_{C(\tau)/S(\tau)}$ which can be expressed as
$
W \cong \bigl\{(z,w)\in \mathbb{C}^2\mid | w| <r \bigr\}/{\sim}$
for a real number~${r>1}$, where $\sim$ is the equivalence relation generated by
\[(z,w)\sim\bigl(z+1,\exp\bigl(p\cdot 2\pi \sqrt{-1}\bigr)\cdot w\bigr)\sim\bigl(z+\tau,\exp\bigl(q\cdot 2\pi \sqrt{-1}\bigr)\cdot w\bigr)\]
for $\tau\in Y$.

Then for $s\in \Delta\backslash \{0\}$ with $\Delta:=\{s\in \mathbb{C}\mid |s|<1\}$, let
\[M_s:=S(\tau)\backslash\bigl\{[(z,w)]\in W\mid |w| \leq \sqrt{|s|}/r\bigr\}\]
and
\[V_s:=\bigl\{[(z,w)]\in W\mid \sqrt{|s|}/r<|w|<\sqrt{|s|}r\bigr\}.\]

The following is to describe some ample line bundles on $S(\tau)$.

Let $E_i$ be the exceptional divisor corresponding to the point $p_i \in Z:= \{p_1, p_2, \dots , p_9 \}$ for~${i \in \{1, 2, \dots , 9 \}}$, respectively. Moreover, let $\pi\colon S(\tau)\rightarrow \operatorname{CP}^2$ be the blow up of $\operatorname{CP}^2$ at nine points in $Z:= \{p_1, p_2, \dots , p_9 \}\subset f_{\tau}(C_0(\tau))$ and $H=(\pi)^* \mathcal{O}_{\operatorname{CP}^2}(1)$ be the pull back bundle of the hyperplane line bundle. Then the Nakai--Moishezon criterion \cite[Appendix~A]{Robin} can be used to prove the following simplified theorem.

\begin{Lemma}[{\cite[Theorem~3]{Tomasz-Halszka}}] \label{Lemma2.5}
Let $\pi\colon S(\tau)\rightarrow$ \emph{$\operatorname{CP}^2$} be the blowing up of \emph{$\operatorname{CP}^2$} at nine points defined above. For the line bundles of the form $L=d\cdot H-k\cdot\sum _{i=1}^9 E_i$ with $k\geq2$ and $d\geq3k+1$ being integers, if \smash{$9\leq \frac{d^2}{k^2}-1$}, then the line bundles are ample.
\end{Lemma}

Now let $C^{\pm}_0(\tau) \subset S^{\pm}_0=\operatorname{CP}^2 $ be two copies of $f_{\tau}(C_0(\tau))$. Then the map
\[\ell_{\operatorname{CP}^2}([a_1:a_2:a_3])=[a_1:-a_2:a_3] \qquad \text{for} \quad [a_1:a_2:a_3]\in \operatorname{CP}^2\]
is a holomorphic involution on $\operatorname{CP}^2$ and induces a holomorphic involution on $f_{\tau}(C_0(\tau))$ as well as \smash{$C^\pm_0(\tau)$} (i.e., \smash{$(\ell_{\operatorname{CP}^2})^{-1}=\ell_{\operatorname{CP}^2}$}). So there is a biholomorphic map \smash{$\ell_{S^{+}_0}$} induced by \smash{$\ell_{\operatorname{CP}^2}$} from~$S^{+}_0$ to $S^{-}_0$. Since $C^\pm_0(\tau)$ are two copies of $f_{\tau}(C_0(\tau))$, there exists an identity map $g_\tau$ from~$C_0^{+}(\tau)$ to~$C_0^{-}(\tau)$. Let \smash{$Z^{\pm}=\bigl\{p^{\pm}_1, p^{\pm}_2, \dots , p^{\pm}_9\bigr\}$} with \smash{$p^{-}_i=\ell_{S^{+}_0}\bigl(p^{+}_i\bigr)$} for $i \in \{1, 2, \dots , 9 \}$ be the sets collecting the nine blow up points on \smash{$C^\pm_0(\tau)$}, respectively. Moreover, let $p^{+}_0$ and~\smash{$p^{-}_0=\ell_{S^{+}_0}\bigl(p^{+}_0\bigr)$} be inflection points of \smash{$C_0^{\pm}(\tau)$}, respectively. Here, the set \smash{$Z^+ \subset C^+_0(\tau)$} is just a~copy of the set~$Z \subset f_{\tau}(C_0(\tau))$.

In addition, let \smash{$S^{\pm}(\tau)\cong\operatorname{CP}^2 \#9\overline{\operatorname{CP}}^2$} be the blow-ups of $S^{\pm}_0$ at nine points in the sets $Z^{\pm}:=\bigl\{p^{\pm}_1, p^{\pm}_2, \dots , p^{\pm}_9\bigr\}$ on $C_0^{\pm}(\tau)$ and $C^{\pm}(\tau)$ be the strict transforms of $C_0^{\pm}(\tau)$, respectively. Furthermore, let \smash{$\ell_{S^{+}(\tau)}\colon S^{+}(\tau)\rightarrow S^{-}(\tau)$} be the biholomorphic map defined on $S^{+}(\tau)$ induced by the biholomorphic map \smash{$\ell_{S^{+}_0}\colon S^{+}_0\rightarrow S^{-}_0$}.

Then the biholomorphic map $\widetilde{g}_{\tau}\colon C^{+}(\tau)\rightarrow C^{-}(\tau)$ induced by $g_{\tau}$ will map the normal bundle~${N_{C^+(\tau)/S^+(\tau)}\in \operatorname{Pic}^{0}\bigl(C^{+}(\tau)\bigr)}$ to the dual of the normal bundle $N_{C^-(\tau)/S^-(\tau)}\in \operatorname{Pic}^{0}(C^{-}(\tau))$ through $(\widetilde{g}_{\tau})^*N_{C^-(\tau)/S^-(\tau)}\cong (N_{C^+(\tau)/S^+(\tau)})^{-1}$. This can be derived through computing the monodromies of \smash{$N_{C^{\pm}(\tau)/S^{\pm}(\tau)}$}.

Taking use of Lemma~\ref{Lemma2.5}, let $E^{\pm}_i$ be the exceptional divisors corresponding to the points $p^{\pm}_i \in Z^{\pm}=\bigl\{p^{\pm}_1, p^{\pm}_2, \dots , p^{\pm}_9\bigr\}$ for $i \in \{1, 2, \dots , 9 \}$, respectively. Moreover, let $H^{\pm}$ be the pull back bundles of the hyperplane line bundles on $S_0^{\pm}$ respectively with $H^+$ being a copy of $H$ and $(\ell_{S^{+}(\tau)})^*H^-=H^+$. So there are ample line bundles $L^{\pm}$ on $S^{\pm}(\tau)$ such that $\bigl(L^{+}.C^{+}(\tau)\bigr)=(L^{-}.C^{-}(\tau))$ with $(\ell_{S^{+}(\tau)})^*L^{-}=L^{+}$ as follows.

\begin{Theorem} \label{Theorem2.6}
There exist ample line bundles $L^{\pm}$ on $S^{\pm}(\tau)$ such that $\bigl(L^{+}.C^{+}(\tau)\bigr)\!=\!(L^{-}.C^{-}(\tau))$ with $(\ell_{S^{+}(\tau)})^*L^{-}=L^{+}$.
\end{Theorem}

\begin{proof}
Since \smash{$K^{-1}_{S^{+}(\tau)}=\big[C^{+}(\tau)\big]=3H^{+}-\sum _{i=1}^9 E_i^{+}$} and \smash{$K^{-1}_{S^{-}(\tau)}=[C^{-}(\tau)]=3H^{-}-\sum _{i=1}^9 E_i^{-}$}, \smash{$\bigl(L^{+}.C^{+}(\tau)\bigr)=3d^{+}-9k^{+}$} and \smash{$(L^{-}.C^{-}(\tau))=3d^{-}-9k^{-}$}. Therefore, for $L^{+}=d^{+}\cdot H^{+}-k^{+} \cdot \smash{\sum _{i=1}^9 E_i^{+}}$ and \smash{$L^{-}=d^{-}\cdot H^{-}-k^{-}\cdot\sum _{i=1}^9 E_i^{-}$} with $k^{\pm}\geq2$, $d^{+} \geq 3k^{+}+1 $ and $d^{-} \geq 3k^{-}+1 $ being integers, $\bigl(L^{+}.C^{+}(\tau)\bigr)=(L^{-}.C^{-}(\tau))$ if and only if $d^{+}-3k^{+}=d^{-}-3k^{-}$.

Then for $d^{+}=d^{-}$ and $k^{+}=k^{-}$ satisfying \smash{$10\leq \frac{(d^{+})^2}{(k^{+})^2}\leq \frac{(d^{-})^2}{(k^{-})^2}$} with $k^{+}=k^{-}\geq 2$ and $d^{+}=d^{-}\geq 3k^{+}+1=3k^{-}+1$ being integers,
\[
L^{+}=d^{+}\cdot H^{+}-k^{+}\cdot\sum _{i=1}^9 E_i^{+}\qquad \text{and}\qquad L^{-}=d^{-}\cdot H^{-}-k^{-}\cdot\sum _{i=1}^9 E_i^{-}\]
 are ample line bundles on $S^{+}(\tau)$ and $S^{-}(\tau)$ respectively such that $\bigl(L^{+}.C^{+}(\tau)\bigr)=(L^{-}.C^{-}(\tau))$. In addition, $(\ell_{S^{+}(\tau)})^*L^{-}=L^{+}$. So Theorem~\ref{Theorem2.6} is proved.
\end{proof}


In fact, for an arbitrary ample line bundle $L^+$ on $S^+(\tau)$ and a line bundle $L^-$ on $S^-(\tau)$ defined through $(\ell_{S^{+}(\tau)})^*L^{-}=L^{+}$, it is not so hard to prove that $\bigl(L^{+}.C^{+}(\tau)\bigr)=(L^{-}.C^{-}(\tau))$. So the construction in this paper can be generalized.

Now let $\alpha_{\pm}$ and $\beta_{\pm}$ be smooth loops on $C^{\pm}(\tau)$ corresponding to the line segments $[0,1]$ and $[0,\tau]$ on $C_0(\tau)$, respectively. So \smash{${\rm e}^{{\pm}p\cdot2\pi\sqrt{-1}}$} and \smash{${\rm e}^{{\pm}q\cdot2\pi\sqrt{-1}}$} are the monodromies of $N_{C^{\pm}(\tau)/S^{\pm}(\tau)}$ along the loops $\alpha_{\pm}$ and $\beta_{\pm}$, respectively with $(p,q) \in \mathbb{R}^2$ being a Diophantine number pair. Taking use of Theorem~\ref{Theorem2.4}, the following statement is true.

There are tubular neighborhoods $W^{\pm}$ of $C^{\pm}(\tau)$ biholomorphic to neighborhoods of the zero section in \smash{$N_{C^{\pm}(\tau)/S^{\pm}(\tau)}$} such that $W^+$ is just a copy of $W$ and \smash{$W^- =\ell_{S^+(\tau)}\bigl(W^+\bigr)$} for $\tau\in Y$.

Then for arbitrary ample line bundles $L^{\pm}\rightarrow S^{\pm}(\tau)$ satisfying $\bigl(L^{+}.C^{+}(\tau)\bigr)=(L^{-}.C^{-}(\tau))$ with $(\ell_{S^{+}(\tau)})^*L^{-}=L^{+}$, let
\[p^{\pm}=3f^{-1}_{\tau}\bigl(p_0^{\pm}\bigr)H^{\pm}-\sum_{i=1}^9f^{-1}_{\tau}\bigl(p_i^{\pm}\bigr)E^{\pm}_i\]
and $b_0=\bigl(L^{+}.C^{+}(\tau)\bigr)=(L^{-}.C^{-}(\tau))$.Then there exists a unique $\xi=\frac{1}{b_0}\bigl((p^-.L^-)-\bigl(p^+.L^+\bigr)\bigr)\in\mathbb{C}$ up to modulo $\langle1,\tau\rangle$ such that $g_{\xi}=\ell_{\xi}\circ\widetilde{g}_{\tau}$ with $\ell_{\xi}$ being a translation induced from $\mathbb{C}\ni z\rightarrow z+\xi\in\mathbb{C}$ gives the following gluing construction \cite[Proposition~5.1]{Takayuki-Takato2}. In our case, $\xi =0$ and therefore $g_{\xi}$ coincides with $\widetilde{g}_{\tau}$.

For $s\in \Delta \backslash \{0\}$ with $\Delta:=\{s\in \mathbb{C}\mid |s|<1\}$, let $M_s^+\subset S^+(\tau)$ be a copy of $M_s$, $M^-_s =\ell_{S^+(\tau)}\bigl(M^+_s\bigr) \subset S^-(\tau)$, $V^+_s \subset S^+(\tau)$ be a copy of $V_s$ and $V^-_s=\ell_{S^+(\tau)}\bigl(V^+_s\bigr) \subset S^-(\tau)$ with
$f_s\colon V^{+}_s \rightarrow V^{-}_s$ defined by
\[f_s\bigl(\big[\bigl(z^+,w^+\bigr)\big]\bigr)=\big[\bigl(g_{\xi}\bigl(z^+\bigr),s/w^+\bigr)\big] \qquad \text{for} \quad \big[\bigl(z^+,w^+\bigr)\big]\in V^+_s.\]

The following is the definition for the symmetric K3 surfaces here.

\begin{Definition} \label{Definition2.7}
For a K3 surface $\mathcal{K}$, if there exists an ample line bundle $L_{\mathcal{K}}$ together with a~non-trivial holomorphic involution $f_{\mathcal{K}}$ defined on $\mathcal{K}$ such that $f_{\mathcal{K}}^*L_{\mathcal{K}}=L_{\mathcal{K}}$, then $\mathcal{K}$ is said to be symmetric with respect to $L_{\mathcal{K}}$.
\end{Definition}

So the gluing construction of symmetric K3 surfaces is as follows.

\begin{Theorem} \label{Theorem2.8}
There exists a~sufficiently small $\varepsilon_0 >0$ such that for any $s \in \Delta \setminus \{ 0\}$ and ${0<|s|<\varepsilon_0}$, identifying $V^{+}_s$ and $V^{-}_s$ through the biholomorphic map $f_s$ defined above, the two open complements $M^{\pm}_s$ of the closures of the tubular neighborhoods $W^{\pm}$ of $C^{\pm}(\tau)$ are glued to be a~symmetric projective K3 surface $X_s$.
\end{Theorem}

\begin{proof}
Let $\ell_{S^+(\tau)\setminus C^+(\tau)}$ be defined as $\ell_{S^+(\tau)\setminus C^+(\tau)}=\ell_{S^+(\tau)}|_{S^+(\tau)\setminus C^+(\tau)}$ for $\tau \in Y$. Since $\ell_{S^+(\tau)}\colon S^+(\tau)\rightarrow S^-(\tau)$ is a biholomorphic map with $\ell_{S^+(\tau)}|_{C^{+}(\tau)}\colon C^+(\tau)\rightarrow C^-(\tau)$ being biholomorphic, taking use of \cite[Proposition~2.1]{Takayuki-Takato1}, there are global holomorphic non-vanishing 2-forms~$\eta^{\pm}_s$ defined on $S^{\pm}(\tau)\setminus C^{\pm}(\tau)$ with \smash{$\ell^{*}_{S^{+}(\tau)\setminus C^{+}(\tau)}(\eta^{-}_s)=\eta^{+}_s$} such that the restrictions of~$\eta^{\pm}_s$ on two open complements $M^{\pm}_s$ of the closures of the tubular neighborhoods of $C^{\pm}(\tau)$ can be glued through $f_s$ to be a global holomorphic non-vanishing 2-form $\sigma_s$ on $X_s$. In addition,
\[\sigma_s|_{V^+_s}=a\cdot \frac{{\rm d}z^+ \wedge {\rm d}w^+}{w^+}\]
for $a \in \mathbb{C} \setminus \{ 0 \}$ . So $\sigma_s$ can be assumed to be normalized in this paper. That is to say, $a=1$. Then $X_s$ is a K3 surface.

\begin{Lemma}[{\cite[Theorem~1.2]{Takayuki-Takato2}}]\label{Lemma2.9} Let $L_s$ be the holomorphic line bundle on $X_s$ derived from $L^{\pm}$ through the gluing construction above. Then $\exists$ sufficiently small $\varepsilon_0 >0$ such that $L_s$ is ample for any $s \in \Delta \setminus \{ 0\} $ and $0<|s|<\varepsilon_0$.
\end{Lemma}

Taking use of Lemma~\ref{Lemma2.9}, for any $s \in \Delta \setminus \{ 0\} $ and $0<|s|<\varepsilon_0$, $L_s$ is an ample line bundle on $X_s$. So $X_s$ is projective for $s \in \Delta \setminus \{ 0\} $ and $0<|s|<\varepsilon_0$.

Let $F_{X_s}\colon X_s \rightarrow X_s$ be defined by
\[F_{X_s}(x)=\begin{cases}
\ell_{M^+_s}(x), &x \in M^{+}_s,\\
\ell^{-1}_{M^+_s}(x), &x \in M^{-}_s
\end{cases}\]
with \smash{$\ell_{M^{+}_s}$} being the restriction of $\ell_{S^{+}(\tau)}$ on $M^{+}_s$. Then $F_{X_s}$ is a non-trivial holomorphic involution on the K3 surface $X_s$. In addition, since $(\ell_{S^{+}(\tau)})^*L^{-}=L^{+}$ , $F^{*}_{X_s}L_s=L_s$ for $s \in \Delta \setminus \{ 0\} $ and~${0<|s|<\varepsilon_0}$. So $X_s$ is a symmetric projective K3 surface for $s \in \Delta \setminus \{ 0\} $ and $0<|s|<\varepsilon_0$. Then Theorem~\ref{Theorem2.8} is proved.
\end{proof}


Here, the existence of $\varepsilon_0$ will also be confirmed in Section~\ref{sec3} during the process of the construction of complete K\"ahler metrics on the quasi-projective varieties $S(\tau)\setminus C(\tau)$ for $\tau \in Y$.

\subsection{Toroidal groups and tubular neighborhoods}

The following is to introduce the corresponding toroidal groups $U_0=\mathbb{C}^2/\Lambda_0$ with
$\Lambda_0=\big\langle\big(\begin{smallmatrix}0 \\1\end{smallmatrix}\big),\allowbreak\big(\begin{smallmatrix}1 \\p\end{smallmatrix}\big),\big(\begin{smallmatrix}\tau \\q\end{smallmatrix}\big)\big\rangle$
provided in \cite{Takayuki-Takato2} for the Diophantine pair $(p,q)$ and show that $U_0=\mathbb{C}^2/\Lambda_0$ is a~quasi-abelian variety.

\begin{Definition}[{\cite[Definition~1.1.1]{Yukitaka-Klaus}}] \label{Definition2.10}
 Let $\Lambda$ be a lattice which is a discrete subgroup of $\mathbb{C}^n$ for $n\in \mathbb{N}$. Then $\mathbb{C}^n/\Lambda$ is called a toroidal group if there does not exist any non-constant holomorphic function on $\mathbb{C}^n/\Lambda$.
\end{Definition}

The toroidal group is a topological group. Moreover, it is an abelian complex Lie group~\cite{Yukitaka-Klaus}.

\begin{Lemma}[{\cite[Theorem~1.1.4]{Yukitaka-Klaus}}] \label{Lemma2.11}
 For $\mathbb{C}^2/\Lambda$ with $\Lambda$ a discrete subgroup of $\mathbb{C}^2$, $\mathbb{C}^2/\Lambda$ is a~toroidal group if and only if there does not exist any~${\sigma \in \mathbb{C}^2 \backslash \{ 0 \}}$ such that the scalar product~${\langle \sigma,\lambda \rangle \in \mathbb{Z}}$ is integral for all $ \lambda \in \Lambda $.
\end{Lemma}

Taking use of Lemma~\ref{Lemma2.11}, the following lemma can be derived. Here, let $\mathbb{Q}$ denote the set consisting of rational numbers.

\begin{Lemma} \label{Lemma2.12}
For any real number pair $(p,q)$ satisfying the Diophantine condition and complex number $\tau \in Y$, $U_0=\mathbb{C}^2/\Lambda_0$ with $\Lambda_0=\bigl\langle\bigl(\begin{smallmatrix}0 \\1\end{smallmatrix}\bigr),\bigl(\begin{smallmatrix}1 \\p\end{smallmatrix}\bigr),\bigl(\begin{smallmatrix}\tau \\q\end{smallmatrix}\bigr)\bigr\rangle$ is a toroidal group.
\end{Lemma}

\begin{proof}
For $\sigma =\bigl(\begin{smallmatrix}\sigma_1\\\sigma_2\end{smallmatrix}\bigr) \in \mathbb{C}^2 \backslash \{ 0 \}$,
\begin{gather*}
\left\langle \sigma,\begin{pmatrix}0 \\1\end{pmatrix} \right\rangle =\left\langle \begin{pmatrix}\sigma_1\\
\sigma_2\end{pmatrix}, \begin{pmatrix}0 \\1\end{pmatrix} \right\rangle =\sigma_2, \qquad \left\langle \sigma,\begin{pmatrix}1 \\p\end{pmatrix} \right\rangle=\left\langle \begin{pmatrix}\sigma_1\\
\sigma_2\end{pmatrix}, \begin{pmatrix} 1 \\p \end{pmatrix} \right\rangle =\sigma_1+p\sigma_2
\end{gather*}
and $\bigl\langle \sigma\bigl(\begin{smallmatrix}\tau \\q\end{smallmatrix}\bigr) \bigr\rangle =\bigl\langle \bigl(\begin{smallmatrix}\sigma_1\\
\sigma_2\end{smallmatrix}\bigr), \bigl(\begin{smallmatrix} \tau \\ q \end{smallmatrix}\bigr) \bigr\rangle =\sigma_1 \tau+q\sigma_2$.

If $q \in \mathbb{Q}$, since $(p,q) \in \mathbb{R}^2$ satisfies the Diophantine condition, $p \notin \mathbb{Q}$. Under the condition that $\sigma_2 \in \mathbb{Z}$ and $\sigma_1 \tau + q\sigma_2\in \mathbb{Z}$, then $\sigma_1=0$ or $\sigma_1 \in \mathbb{C} \setminus \mathbb{R}$. So $\sigma_1 + p\sigma_2\notin \mathbb{Z}$.

If $q \notin \mathbb{Q}$, under the condition that $\sigma_2 \in \mathbb{Z}$ and $\sigma_1 \tau + q\sigma_2\in \mathbb{Z}$, then $\sigma_1 \in \mathbb{C} \setminus \mathbb{R}$. So $\sigma_1 + p\sigma_2\notin \mathbb{Z}$.

Therefore, there does not exist any $\sigma \in \mathbb{C}^2 \backslash \{ 0 \}$ such that the scalar product $\langle \sigma,\lambda \rangle \in \mathbb{Z}$ is integral for all $ \lambda \in \Lambda_0 $.

Taking use of Lemma~\ref{Lemma2.11}, $U_0=\mathbb{C}^2/\Lambda_0$ with $\Lambda_0=\bigl\langle\bigl(\begin{smallmatrix}0 \\1\end{smallmatrix}\bigr),\bigl(\begin{smallmatrix}1 \\p\end{smallmatrix}\bigr),\bigl(\begin{smallmatrix}\tau \\q\end{smallmatrix}\bigr)\bigr\rangle$ is a toroidal group for any $(p,q) \in \mathbb{R}^2$ satisfying the Diophantine condition. So Lemma~\ref{Lemma2.12} is proved.
\end{proof}


\begin{Theorem}[{\cite[Proposition~2.4]{Takayuki-Takato2}}] \label{Theorem2.13}
 For $(p,q) \in \mathbb{R}^2$ satisfying the Diophantine condition and $\tau \in Y$, every topologically trivial line bundle on the toroidal group defined as $U_0=\mathbb{C}^2/\Lambda_0$ with $\Lambda_0=\bigl\langle\bigl(\begin{smallmatrix}0 \\1\end{smallmatrix}\bigr),\bigl(\begin{smallmatrix}1 \\p\end{smallmatrix}\bigr),\bigl(\begin{smallmatrix}\tau \\q\end{smallmatrix}\bigr)\bigr\rangle$ is homogeneous.
\end{Theorem}

\begin{proof}
Taking use of the result from \cite{Christian}, there are nine equivalent conditions for the special class of toroidal groups with homogeneous topologically trivial line bundles \cite[Theorem]{Christian}. The following lemma is to show the rewritten condition 9.

\begin{Lemma}[{\cite{Christian}}] \label{Lemma2.14}
 For $(p,q) \in \mathbb{R}^2$ satisfying the Diophantine condition and $\tau \in Y$, let $D_0=\bigl(\begin{matrix}p& q\end{matrix}\bigr)$. Then every topologically trivial line bundle on the toroidal group defined as $U_0=\mathbb{C}^2/\Lambda_0$ with $\Lambda_0=\bigl\langle\bigl(\begin{smallmatrix}0 \\1\end{smallmatrix}\bigr),\bigl(\begin{smallmatrix}1 \\p\end{smallmatrix}\bigr),\bigl(\begin{smallmatrix}\tau \\q\end{smallmatrix}\bigr)\bigr\rangle$ is homogeneous if and only if there exist $ c>0$ and $a \geq 0$ such that~${\|\sigma^t \cdot D_0+ \delta^t\| \geq c{\rm e}^{-a|\sigma|}}$
for all $\sigma \in \mathbb{Z} \setminus \{ 0 \}$ and $\delta \in \mathbb{Z}^2$.
\end{Lemma}

Since $(p,q)\in \mathbb{R}^2$ satisfying the Diophantine condition and $\tau \in Y$, there exist $\vartheta >0$ and~${A>0}$ such that
\[
\underset{\mu_0 ,\nu_0 \in \mathbb{Z}}{\min } \big| n\bigl(p+q\sqrt{-1} \bigr)-\bigl(\mu_0 +\nu_0\sqrt{-1} \bigr)\big| \geq A\cdot n^{-\vartheta }
\]
for any $n \in \mathbb{N}$.

That is to say,
\[\| \sigma^t \cdot D_0+ \delta^t\|=\| \sigma(p\quad q)+ \delta^t\| \geq A\cdot {\rm e}^{-\vartheta|\sigma|}\]
for all $\sigma \in \mathbb{Z} \backslash \{ 0 \}$ and $\delta \in \mathbb{Z}^2$. Therefore, the nine equivalent conditions for the special class of toroidal groups \cite[Theorem]{Christian} are all satisfied here. That is to say, every topologically trivial line bundle on the toroidal group defined as $U_0=\mathbb{C}^2/\Lambda_0$ with $\Lambda_0=\bigl\langle\bigl(\begin{smallmatrix}0 \\1\end{smallmatrix}\bigr),\bigl(\begin{smallmatrix}1 \\p\end{smallmatrix}\bigr),\bigl(\begin{smallmatrix}\tau \\q\end{smallmatrix}\bigr)\bigr\rangle$ is homogeneous. So Theorem~\ref{Theorem2.13} is proved.
\end{proof}


With the above result, from \cite{Christian}, the complex vector space $H^1(U_0,\mathcal{O})$ is finite-dimensional. The following is to show that $U_0=\mathbb{C}^2/\Lambda_0$ defined above are quasi-abelian varieties.

Let
\[\mathbb{R}_{\Lambda_{0}}:=\left \{ x_1 \cdot \begin{pmatrix}0 \\1\end{pmatrix}+x_2 \cdot \begin{pmatrix}1 \\p\end{pmatrix} +x_3 \cdot \begin{pmatrix}\tau \\q\end{pmatrix}\mid x_1\in \mathbb{R},\,x_2\in \mathbb{R},\,x_3\in \mathbb{R}\right \}\]
be the $\mathbb{R}$-span of $\Lambda_0$. Let
\[M\mathbb{C}_{\Lambda_0}:=\mathbb{R}_{\Lambda_0}\cap\sqrt{-1}\mathbb{R}_{\Lambda_0}=\left \{ \begin{pmatrix}x \\0\end{pmatrix}\mid x\in \mathbb{C}\right \}\]
be the maximal $\mathbb{C}$-linear subspace of $\mathbb{R}_{\Lambda_{0}}$.

\begin{Definition}[{\cite[Definition~3.1.6]{Yukitaka-Klaus}}]\label{Definition2.15}
 An ample Riemann form $H_0$ for the discrete subgroup~${\Lambda \subset \mathbb{C}^n}$ of complex rank n is a Hermitian form $H_0$ on $\mathbb{C}^n$ such that
\begin{itemize}\itemsep=0pt
\item[$(1)$]$\operatorname{ Im}H_0|_{\Lambda\times\Lambda}$ is $\mathbb{Z}$-valued,
\item[$(2)$]$H_0$ is positive definite on the maximal $\mathbb{C}$-linear subspace $M\mathbb{C}_\Lambda$ of $\mathbb{R}_\Lambda$.
\end{itemize}
Here, $\mathbb{R}_\Lambda$ is the $\mathbb{R}$-span of $\Lambda$. $M\mathbb{C}_\Lambda$ is the maximal $\mathbb{C}$-linear subspace of $\mathbb{R}_\Lambda$.
\end{Definition}

\begin{Definition}[{\cite[Definition~3.1.6]{Yukitaka-Klaus}}]\label{Definition2.16}
 A quasi-abelian variety is a toroidal group $\mathbb{C}^n/\Lambda$ with an ample Riemann form for the lattice $\Lambda$.
\end{Definition}

Then the following is true.

\begin{Theorem}\label{Theorem2.17}
For $(p,q) \in \mathbb{R}^2$ satisfying the Diophantine condition and $\tau \in Y$, $U_0={\mathbb{C}^2_{(z,\eta)}}/{\Lambda_0}$ with \smash{$\Lambda_0=\bigl\langle\bigl(\begin{smallmatrix}0 \\1\end{smallmatrix}\bigr),\bigl(\begin{smallmatrix}1 \\p\end{smallmatrix}\bigr),\bigl(\begin{smallmatrix}\tau \\q\end{smallmatrix}\bigr)\bigr\rangle$} is a quasi-abelian variety.
\end{Theorem}

\begin{proof}
Let
\[
G=\left(\begin{matrix}\frac{1}{\operatorname{\operatorname{Im}}\tau} & 0\\0 & 0\end{matrix}\right),
\]
 then~$G$ defines a hermitian form on~$\mathbb{C}^2$. $\operatorname{ Im}G|_{\Lambda_0\times \Lambda_0}$ is $\mathbb{Z}$-valued and~$G$ is positive definite on the maximal $\mathbb{C}$-linear subspace $M\mathbb{C}_{\Lambda_0}$ of $\mathbb{R}_{\Lambda_{0}}$. So
 \[
 G=\left(\begin{matrix}\frac{1}{\operatorname{\operatorname{Im}}\tau} & 0\\0 & 0\end{matrix}\right)
 \]
 is an ample Riemann form for the lattice $\Lambda_0$. Therefore, according to Definition~\ref{Definition2.16}, $U_0=\mathbb{C}^2/\Lambda_0$ is a quasi-abelian variety. Then Theorem~\ref{Theorem2.17} is proved.
\end{proof}

Since $M\mathbb{C}_{\Lambda_0}:=\mathbb{R}_{\Lambda_0}\cap\sqrt{-1}\mathbb{R}_{\Lambda_0}=\big \{ \bigl(\begin{smallmatrix}x \\0\end{smallmatrix}\bigr)\mid x\in \mathbb{C}\big \}$ is the maximal $\mathbb{C}$-linear subspace of~$\mathbb{R}_{\Lambda_0}$ and $M\mathbb{C}_{\Lambda_0}$ has the complex dimension 1, the toroidal group $U_0=\mathbb{C}^2/\Lambda_0$ is of type~1 \cite[Defi\-nition~1.2.1]{Yukitaka-Klaus}.

For \smash{$G=\bigl(\begin{smallmatrix}\frac{1}{\operatorname{\operatorname{Im}}\tau} & 0\\0 & 0\end{smallmatrix}\bigr)$}, $\operatorname{\operatorname{Im}}G$ has rank 2 on $\mathbb{R}_{\Lambda_0}$. Moreover, since the toroidal group $U_0=\mathbb{C}^2/\Lambda_0$ is of type 1, the ample Riemann form \smash{$G=\bigl(\begin{smallmatrix}\frac{1}{\operatorname{\operatorname{Im}}\tau} & 0\\0 & 0\end{smallmatrix}\bigr)$} for $\Lambda_0$ is said to be of kind $\frac{2-2\times1}{2}=0$ (see~\cite[Definition~3.1.12]{Yukitaka-Klaus}).

The following theorem will give a fibration of the toroidal group $U_0=\mathbb{C}^2/\Lambda_0$.

\begin{Theorem}[{\cite[Theorem~3.1.16]{Yukitaka-Klaus}}]\label{Theorem2.18}
 For a toroidal group $\mathbb{C}^n/\Lambda$ of type q with $\Lambda$ being a lattice, $\mathbb{C}^n/\Lambda$ is a quasi-abelian variety with an ample Riemann form of kind k if and only if $\mathbb{C}^n/\Lambda$ has a~maximal closed Stein subgroup $K \cong \mathbb{C}^k\times \mathbb{C}^{*m}$ with $2k+m=n-q$ and $(\mathbb{C}^n/\Lambda)/K$ is an abelian variety of dimension $q+k$.
\end{Theorem}

Taking use of Theorem~\ref{Theorem2.18}, the toroidal group $U_0=\mathbb{C}^2/\Lambda_0$ has a maximal closed Stein subgroup $K \cong \mathbb{C}^*$ and $U_0/K$ is an abelian variety of dimension 1.

Furthermore, from \cite{Takayuki-Takato2}, through the map $f_{U_0}\colon U_0=\mathbb{C}^2_{(z,\eta)}/{\Lambda_0}\rightarrow V_{0,\infty}$ defined by sending
\[
[(z,\eta)]\in U_0=\mathbb{C}^2_{(z,\eta)}/{\Lambda_0} \qquad \text{to} \quad \big[\bigl(z,{\rm e}^{2\pi \sqrt{-1}\eta}\bigr)\big]\in V_{0,\infty},
\]
$U_0=\mathbb{C}^2_{(z,\eta)}/{\Lambda_0}\cong V_{0,\infty}$ with $V_{0,\infty}$ defined by
$V_{b_1,b_2}:=\big \{ (z,w)\in \mathbb{C}^2\mid b_1<|w|<b_2 \big \}/{\sim}$
for~${0\leq b_1<b_2\leq\infty}$, where $\sim$ is the equivalence relation generated by
\[(z,w)\sim\big(z+1,\exp\big(p\cdot 2\pi \sqrt{-1}\big)\cdot w)\sim\big(z+\tau,\exp\big(q\cdot 2\pi \sqrt{-1}\big)\cdot w\big)\]
with $\tau \in Y$.

From the definitions above, $V_{0,r}\cup C(\tau)=W$. Let $\pi_W\colon W \rightarrow C(\tau)$ be the natural projection defined as $\pi_W([(z,w)])=[z]$ for $[(z,w)]\in W$. The next step is to give a theorem for the line bundles on the tubular neighborhoods defined above.

\begin{Theorem}[{\cite[Proposition~3.3]{Takayuki-Takato2}}]\label{Theorem2.19}
 For any $L \in\operatorname{ Pic}(W)$, $L$ is the pull back line bundle~$(\pi_W)^*(L|_{C(\tau)})$ of the restricted line bundle $L|_{C(\tau)}$ on $C(\tau)$. Here, the assumption is as above.
\end{Theorem}

This theorem is proved in \cite{Takayuki-Takato2} through the toroidal group defined above.


\section[A deformation family of complete non-compact K\"ahler manifolds]{A deformation family of complete non-compact\\ K\"ahler manifolds}\label{sec3}

\subsection[A deformation family of quasi-projective varieties from CP\string^2 \#9CP\string^2]{A deformation family of quasi-projective varieties from $\boldsymbol{\operatorname{CP}^2 \#9 \overline{\operatorname{CP}}^2}$}\label{sec3.1}

The following is to prove Theorem~\ref{Theorem1.1}. First of all, the definition for a complex analytic family of compact complex manifolds is as follows.

\begin{Definition}[{\cite[Definition~2.8]{Kunihiko}}]\label{Definition3.1}
 For a collection of compact complex manifolds $ {\{M_t \mid t \!\in\! B\}}$ with $B$ being a domain of $\mathbb{C}^m$ and $m\in \mathbb{N} $, $ \{M_t \mid t \in B\}$ is called a complex analytic family of compact complex manifolds if there is a complex manifold $\mathcal{M}$ and a holomorphic map $\pi$ from~$\mathcal{M}$ onto $B$ satisfying the following conditions:
\begin{itemize}\itemsep=0pt
\item[$(1)$] The rank of the Jacobian matrix of $\pi$ is equal to $m$ at every point of $\mathcal{M}$.
\item[$(2)$] For each $t \in B$, $\pi^{-1}(t)=M_t$ is a compact complex submanifold of $\mathcal{M}$.
\end{itemize}
\end{Definition}

Let $(\mathcal{M},B,\pi)$ denote the complex analytic family defined above. Here, a complex analytic family of \smash{$\operatorname{CP}^2 \#9 \overline{\operatorname{CP}}^2$} over a 9-dimensional complex manifold will be presented before constructing a deformation family of quasi-projective varieties.

However, a complex analytic family of elliptic curves over a 1-dimensional complex manifold should be presented at first. Here, for $\tau_0\in Y$, let $U \subset Y$ be a sufficiently small disc neighborhood of $\tau_0$ in $Y$ and $\pi_{\operatorname{CP}^2 \times U}\colon \operatorname{CP}^2\times U\rightarrow U$ be the natural projection map. Let $C_0(\tau)$ be an elliptic curve embedded in
\smash{$\pi_{\operatorname{CP}^2\times U}^{-1}(\tau)=\operatorname{CP}^2 \times\{\tau\}$} as $f_{\tau}(C_0(\tau))\times \{ \tau \}$
for $\tau \in U$ and
\[\widetilde{\mathcal{S}}=\bigl\{ (x, \tau) \in \operatorname{CP}^2 \times U\mid x \in f_{\tau}(C_0(\tau)) \bigr\}.
\]
 Here, $C_0( \tau_0)=C_0$. Then the following lemma is true.

\begin{Lemma}\label{Lemma3.2}
\emph{$\bigl(\widetilde{\mathcal{S}},U,\pi_{\operatorname{CP}^2 \times U}|_{\widetilde{\mathcal{S}}}\bigr)$} is a complex analytic family of elliptic curves over the $1$-di\-men\-sion\-al complex manifold $U$ and \emph{$\bigl(\operatorname{CP}^2\times U, U,\pi_{\operatorname{CP}^2\times U}\bigr)$} is a complex analytic family of \emph{ $\operatorname{CP}^2$} over the $1$-dimensional complex manifold $U$.
\end{Lemma}

The complete proof of Lemma~\ref{Lemma3.2} is in Appendix~\ref{appendixA}.

Now let $U \subset Y$ be a sufficiently small disc neighborhood of $\tau_0$ and $U_\nu$ be sufficiently small disc neighborhoods of $q_\nu$ in the universal cover $\mathbb{C}$ of $ \mathbb{C}/ \langle 1,\tau \rangle=C_0({\tau})$ for $\nu \in \{1,2,3,4,5,6,7,8 \}$. Here, $C_0(\tau_0)=C_0$. Let $f_{\widetilde{\tau},\nu}\colon U_{\nu}\!\rightarrow U_{_\nu}/\langle 1, \widetilde{\tau}\rangle$ be the natural projection for $\nu \in \{1,2,3,4,5,6,7,8 \}$ and $\widetilde{\tau}\in U$.

Let $(p,q)$ still be a Diophantine number pair. For $\hat{t}=(\widetilde{\tau},\hat{p}_1,\dots ,\hat{p}_8)\in T =U \times U_1 \times \dots \times U_8$, let $f_9\bigl(\hat{t}\bigr) \in C_0(\widetilde{\tau})$ be the point fixed by the equation
\[9p_0-\sum _{j=1}^8f_{\widetilde{\tau},j}(\hat{p}_j)-f_9\bigl(\hat{t}\bigr)=q-p\cdot \widetilde{\tau} \quad \mod \langle 1, \widetilde{\tau} \rangle \]
with $p_0$ being an inflection point \cite{Takayuki-Takato2}. In addition, let $\pi\colon\operatorname{CP}^2 \times T \rightarrow T$ be the natural projection and let $\mathcal{S}_0=\widetilde{\mathcal{S}} \times U_1 \times \dots \times U_8$. Then $\bigl(\operatorname{CP}^2 \times T,T, \pi\bigr)$ is a complex analytic family of $\operatorname{CP}^2$.

Through blowing up nine points $ \{ f_{\widetilde{\tau}}(p_1),f_{\widetilde{\tau}}(p_2),\dots ,f_{\widetilde{\tau}}(p_9) \} \times \{ (\widetilde{\tau},\hat{p}_1,\dots ,\hat{p}_8) \}$ on
\[\pi|^{-1}_{\mathcal{S}_0}\bigl(\hat{t}\bigr)=\pi_{\operatorname{CP}^2 \times U}|^{-1}_{\widetilde{S}}(\widetilde{\tau})\times ] \{ (\hat{p}_1,\dots ,\hat{p}_8) ] \}
\]
 with
 \[
 \big \{ p_1=f_{\widetilde{\tau},1}(\hat{p}_1),\dots ,p_8=f_{\widetilde{\tau},8}(\hat{p}_8),p_9=f_9\bigl(\hat{t}\bigr) \big \}\subset \mathbb{C}/ \langle 1, \widetilde{\tau} \rangle
 \]
 for each ${\hat{t}=(\widetilde{\tau},\hat{p}_1,\dots ,\hat{p}_8)\in T}$, a~11-dimensional complex manifold $\mathcal{F}$ can be derived from $\operatorname{CP}^2 \times T$. Moreover, let $\widetilde{\pi}$ be defined as a map from $\mathcal{F}$ to $T$ such that \smash{$(\widetilde{\pi})^{-1}\bigl(\hat{t}\bigr)$} is the blow up of \smash{$\pi^{-1}\bigl(\hat{t}\bigr)$} at the nine points in~$ { \{ f_{\widetilde{\tau}}(p_1),\dots ,f_{\widetilde{\tau}}(p_9) \} \times \{ (\widetilde{\tau},\hat{p}_1,\dots ,\hat{p}_8) \}}$ with
\[
\{ p_1=f_{\widetilde{\tau},1}(\hat{p}_1),\dots , p_8=f_{\widetilde{\tau},8}(\hat{p}_8),p_9=f_9\bigl(\hat{t}\bigr) \}\subset \mathbb{C}/ \langle 1, \widetilde{\tau} \rangle
\]
for each $\hat{t}=(\widetilde{\tau},\hat{p}_1,\dots ,\hat{p}_8)\in T$. Now the normal bundle of the strict transform of the elliptic curve $f_{\widetilde{\tau}}(C_0(\widetilde{\tau}))\times \{ (\widetilde{\tau}, \hat{p}_1,\dots ,\hat{p}_8) \}$ in $(\widetilde{\pi})^{-1}\bigl(\hat{t}\bigr)$ satisfies the Diophantine condition corresponding to $(p,q)$ for $\hat{t}=(\widetilde{\tau},\hat{p}_1,\dots ,\hat{p}_8)\in T$ \cite{Takayuki-Takato2}. Let $\mathcal{S} \subset \mathcal{F}$ be the strict transform of~${\mathcal{S}_0 \subset \operatorname{CP}^2 \times T}$. Then the following lemma is true.

\begin{Lemma}[{\cite{Takayuki-Takato1}}] \label{Lemma3.3}
 Through blowing up nine points on every fiber of \emph{$\bigl(\operatorname{CP}^2 \times T,T, \pi\bigr)$} such that the normal bundle of the strict transform of the elliptic curve mentioned above embedded in each fiber satisfies the Diophantine condition corresponding to the Diophantine number pair $(p,q)$, the complex analytic families \emph{$(\mathcal{F},T, \widetilde{\pi})$} and \emph{$(\mathcal{S},T, \widetilde{\pi}|_\mathcal{S})$} can be constructed.
\end{Lemma}

The complete proof of Lemma~\ref{Lemma3.3} is in Appendix~\ref{appendixA}. Here, $(\mathcal{F},T, \widetilde{\pi})$ is a complex analytic family of \smash{$\operatorname{CP}^2 \#9 \overline{\operatorname{CP}}^2$} over the 9-dimensional complex manifold $T$ and $(\mathcal{S},T, \widetilde{\pi}|_\mathcal{S})$ is a complex analytic family of elliptic curves over the 9-dimensional complex manifold $T$.

Secondly, the following is to give a deformation family of quasi-projective varieties.

Here, the fibers of deformation families defined in Definition~\ref{Definition3.1} should be compact. Since there will be a deformation family of open complex manifolds to be confirmed, the definitions of deformation in \cite{Elizabeth-Francisco} will be used. Then the corresponding definition for the category of smooth manifolds is as follows.

\begin{Definition}[the category of smooth manifolds] \label{Definition3.4}
 The objects for the category of smooth manifolds are the smooth manifolds. For any objects $M$ and $N$, the morphisms from $M$ to $N$ are the maps in the set $C^\infty (M,N)$.
\end{Definition}

So the definition of deformation family from \cite{Elizabeth-Francisco} is as follows.

\begin{Definition}[{\cite[Definition~2.3]{Elizabeth-Francisco}}] \label{Definition3.5}
 A deformation family of a complex manifold $X_1$ is a~holomorphic surjective submersion $\mathcal{X} \rightarrow D_1$, where $D_1$ is a complex disc centred at $0$, satisfying
\begin{itemize}\itemsep=0pt
\item[$(1)$] $\pi^{-1}(0)=X_1$,
\item[$(2)$] $\mathcal{X}$ is locally trivial in the $C^\infty$ category.
\end{itemize}
For any $t\in D_1$, the fiber $\pi^{-1}(t)$ is called a deformation of $X_1$.
\end{Definition}

The discussion about Definition~\ref{Definition3.5} can be found in \cite{Edoardo-Elizabeth-Francisco}. Here, there is a simple corresponding lemma.

\begin{Lemma} \label{Lemma3.6}
For a complex analytic family $(\mathcal{M}_1, B_1, \varpi)$ with $B_1$ being a complex disc centred at $0$, $(\mathcal{M}_1, B_1, \varpi)$ is also a deformation family of the complex manifold $\varpi^{-1}(0)$ satisfying the conditions in Definition {\rm\ref{Definition3.5}}.
\end{Lemma}

\begin{proof}
From the Definition~\ref{Definition3.1}, since the rank of the Jacobian matrix of $\varpi$ is equal to 1 at every point of $\mathcal{M}_1$, $\varpi$ is a holomorphic surjective submersion $\mathcal{M}_1 \rightarrow B_1$. Moreover, $\varpi^{-1}(0)$ is a compact complex submanifold of $\mathcal{M}_1$. In addition, the following lemma is from \cite{Kunihiko}.

\begin{Lemma}[{\cite[Theorem~2.5]{Kunihiko}}] \label{Lemma3.7} $\forall x \in B_1$, $\exists$ a polydisc $U_x \subset B_1$ centered at $x$ with a diffeomorphism $\Psi_x \colon \varpi^{-1}(x) \times U_x \rightarrow \varpi^{-1}(U_x)$ such that $\varpi \circ \Psi_x \colon \varpi^{-1}(x) \times U_x \rightarrow U_x$ is a projection.
\end{Lemma}

Taking use of Lemma~\ref{Lemma3.7}, $\varpi$ is $C^\infty$ locally trivial on $B_1$. That is to say, $\mathcal{M}_1$ is locally trivial in the $C^\infty$ category. Therefore, $(\mathcal{M}_1, B_1, \varpi)$ is a deformation family of the complex manifold~$\varpi^{-1}(0)$ satisfying the conditions in Definition~\ref{Definition3.5}. So for a complex analytic family~$(\mathcal{M}_1, B_1, \varpi)$ with $B_1$ being a complex disc centred at $0$, $(\mathcal{M}_1, B_1, \varpi)$ is also a deformation family of the complex manifold $\varpi^{-1}(0)$ satisfying the conditions in Definition~\ref{Definition3.5}. Then Lemma~\ref{Lemma3.6} is proved.
\end{proof}


In addition, the following is to introduce the definition of the smooth compacitifiable deformation family. The complex model space is defined as follows.

\begin{Definition}[{\cite[Definition~13]{Edoardo-Elizabeth-Francisco}}] \label{Definition3.8}
$\forall f_1,\dots ,f_k \in \mathcal{O}(D)$ with $D \subset \mathbb{C}^n$ being a domain and~${n \in \mathbb{N}}$, then $I_D=\mathcal{O}_Df_1+\mathcal{O}_Df_2+\dots +\mathcal{O}_Df_k$ is an ideal sheaf. Let $X_2:=N(f_1,\dots ,f_k)$ and $\mathcal{O}_{X_2}=(\mathcal{O}_D/I_D)|_{X_2}$. So the $\mathbb{C}$-ringed space $(X_2,\mathcal{O}_{X_2})$ is called a complex model space.
\end{Definition}

The complex space is defined as follows.

\begin{Definition}[{\cite[Definition~14]{Edoardo-Elizabeth-Francisco}}] \label{Definition3.9}
For a $\mathbb{C}$-ringed space $\bigl(\hat{X},\mathcal{O}_{\hat{X}}\bigr)$, if $\hat{X}$ is a Hausdorff space and \smash{$\forall x \in \hat{X}$}, $\exists$ an open neighborhood $U_x$ such that the open $\mathbb{C}$-ringed subspace $(U_x,\mathcal{O}_{U_x})$ is isomorphic to a complex model space, then $\bigl(\hat{X},\mathcal{O}_{\hat{X}}\bigr)$ is called a complex space.
\end{Definition}

So the following is the definition of the smooth compactifiable deformation family.

\begin{Definition}[{\cite[Definition~16]{Edoardo-Elizabeth-Francisco}}] \label{Definition3.10}
For a complex space $X$, $X$ is compactifiable if there is an open embedding $\rho\colon X\rightarrow \widetilde{X}$ such that $\rho(X)=\widetilde{X} \backslash D_2$ with $\widetilde{X}$ being a compact complex space and $D_2$ being a closed analytic subset of $\widetilde{X}$. Then a compactifiable deformation of $X$ is a~sixtuple~$(\mathcal{X},D_2,Y_1,f,o,j)$ such that
\begin{itemize}\itemsep=0pt
\item[$(1)$] $\mathcal{X}$ and $Y_1$ are complex spaces,
\item[$(2)$] $D_2$ is a closed analytic subset of $\mathcal{X}$ and $o$ is a point in $Y_1$,
\item[$(3)$] $f$ is a proper holomorphic map from $\mathcal{X}$ to $Y_1$ such that $f|_{\mathcal{X} \backslash D_2}$ is flat,
\item[$(4)$] $j$ is an open embedding from $X$ to $f^{-1}(o)$ with $j(X)=f^{-1}(o)\backslash D_2$.
\end{itemize}
\end{Definition}

Furthermore, the differentially local triviality can be defined as follows.

\begin{Definition}[{\cite[Definition~21]{Edoardo-Elizabeth-Francisco}}] \label{Definition3.11}
For a holomorphic submersion $g\colon X \rightarrow \hat{Y}$ with $X$ and~$\hat{Y}$ being complex manifolds, $g$ is differentially locally trivial over $\hat{Y}$ if $\forall y \in \hat{Y}$, $\exists$ an open neighborhood $U_y$ such that $g|_{g^{-1}(U_y)}$ is differentially isomorphic to the projection ${g^{-1}(y) \times U_y \rightarrow U_y}$.
\end{Definition}

Then for the compactifiable deformation of $X$ defined above, if $f|_{\mathcal{X} \backslash D_2 }\colon \mathcal{X} \backslash D_2\rightarrow Y_1$ with $\mathcal{X} \backslash D_2$ and $Y_1$ being complex manifolds is differentially locally trivial over $Y_1$, then $(\mathcal{X}, D_2, Y_1,f,\allowbreak o,j)$ is said to be differentially trivial along $Y_1$.

\begin{Definition} \label{Definition3.12}
For a compactifiable complex space $X$ defined above, assuming that $X$ is smooth, if $f|_{\mathcal{X}\backslash D_2}$ is a smooth morphism of complex spaces and each fiber of $f$ is a complex manifold, then $(\mathcal{X}, D_2, Y_1,f, o, j)$ is a smooth compactifiable deformation.
\end{Definition}

Taking use of the notations above, let $t_0 = (\tau_0, q_1,q_2,\dots ,q_8) \in T$, then
\[\{ p_1=f_{\tau_0,1}(q_1),\dots ,p_8=f_{\tau_0,8}(q_8),p_9=f_9(t_0) \} \subset C_0(\tau_0)=C_0.\]
Let \smash{$S \cong \operatorname{CP}^2 \#9\overline{\operatorname{CP}}^2$} be the blow up of $S_0 =\operatorname{CP}^2$ at the nine points in the set
\[ \{ f_{\tau_0}(p_1),\dots ,f_{\tau_0}(p_9) \}\subset f_{\tau_0}(C_0)\]
and $C$ be the strict transform of $f_{\tau_0}(C_0)$. Here, the normal bundle $N_{C/S}$ should satisfy the Diophantine condition corresponding to the Diophantine pair $(p,q)$.

A smooth compactifiable deformation of the quasi-projective variety \smash{$S \backslash C \subset S \cong \operatorname{CP}^2 \#9\overline{\operatorname{CP}}^2$} differentially trivial along $T$ is as follows.

\begin{Theorem} \label{Theorem3.13}
Let $o=(\tau_0,q_1,\dots ,q_8)$. Define $j\colon S \backslash C\rightarrow \widetilde{\pi}^{-1}(o)$ by $j(x)=(x,\tau_0, q_1,q_2,\dots ,q_8)$ for $x\in S\backslash C$. Then $(\mathcal{F},\mathcal{S},T,\widetilde{\pi},o,j)$ is a compactifiable deformation of $S\backslash C$.
\end{Theorem}

\begin{proof}
Firstly, it is to prove that $S\backslash C$ is compactifiable.
$S\backslash C$ is a quasi-projective variety. $(\widetilde{\pi})^{-1}(o)\cong S$ is a compact complex manifold. So $(\widetilde{\pi})^{-1}(o)$ is a complex space. Moreover, $(\widetilde{\pi})^{-1}(o)\cap \mathcal{S}\cong \mathbb{C}/\langle 1, \tau_0 \rangle$ is a closed analytic submanifold of $(\widetilde{\pi})^{-1}(o)$. Since $(\widetilde{\pi})^{-1}(o)\backslash j(S\backslash C)=(\widetilde{\pi})^{-1}(o)\cap \mathcal{S}$ is closed, $j(S\backslash C)$ is an open submanifold of~$(\widetilde{\pi})^{-1}(o)$. Then $j$ is an open embedding with $j(S\backslash C)=(\widetilde{\pi})^{-1}(o)\backslash \mathcal{S}$. So $((\widetilde{\pi})^{-1}(o),\widetilde{\pi}^{-1}(o)\cap \mathcal{S},j)$ is a compactification of $S\backslash C$. That is to say, $S\backslash C$ is compactifiable.

Secondly, it is to prove that $\mathcal{S}$ is a closed analytic subset of $\mathcal{F}$ and $\widetilde{\pi}$ is a proper holomorphic map from $\mathcal{F}$ to $T$ such that \smash{$\widetilde{\pi}|_{\mathcal{F}\backslash \mathcal{S}}$} is flat.

Since $\mathcal{F}$ and $T$ are complex manifolds, $\mathcal{F}$ and $T$ are complex spaces. From above, $\mathcal{S}$ is a~codimension 1 complex submanifold of $\mathcal{F}$. In addition, since the blow up points are all on $\mathcal{S}_0$, it is easy to get an open neighborhood for every point in $\mathcal{F}\backslash \mathcal{S}$. Therefore, $\mathcal{F}\backslash \mathcal{S}$ should be an open subset of $\mathcal{F}$. Then $\mathcal{S}$ is a closed subset of $\mathcal{F}$. So $\mathcal{S}$ is a closed analytic subset of $\mathcal{F}$.

Furthermore, since $(\mathcal{F},T,\widetilde{\pi})$ is the complex analytic family of \smash{$\operatorname{CP}^2 \#9\overline{\operatorname{CP}}^2$}, $\widetilde{\pi}$ is a proper holomorphic map from $\mathcal{F}$ to $T$. So $\widetilde{\pi}|_{\mathcal{F}\backslash \mathcal{S}}$ is also proper.

Since $\widetilde{\pi}|_{\mathcal{F}\backslash \mathcal{S}}$ is a proper submersion from the complex manifold $\mathcal{F}\backslash \mathcal{S}$ to $T$, $\widetilde{\pi}|_{\mathcal{F}\backslash \mathcal{S}}$ is flat.

Finally, it is to prove that $(\mathcal{F},\mathcal{S},T,\widetilde{\pi},o,j)$ is a compactifiable deformation of $S\backslash C$.

From above, $j\colon S\backslash C \rightarrow (\widetilde{\pi})^{-1}(o)$ is an open embedding with $j(S\backslash C)=(\widetilde{\pi})^{-1}(o)\backslash \mathcal{S}$. Therefore, with all the proofs above, $(\mathcal{F},\mathcal{S},T,\widetilde{\pi},o,j)$ is a compactifiable deformation of $S\backslash C$. So Theorem~\ref{Theorem3.13} is proved.
\end{proof}


In addition, the following statement is true.

\begin{Proposition} \label{Proposition3.14}
$(\mathcal{F},\mathcal{S},T,\widetilde{\pi},o,j)$ is a smooth compactifiable deformation differentially trivial along $T$.
\end{Proposition}

\begin{proof}
Since \smash{$\widetilde{\pi}|_{\mathcal{F}\backslash \mathcal{S}}$} is a proper submersion from the complex manifold $\mathcal{F}\backslash \mathcal{S}$ to the complex manifold $T$ and \smash{$(\widetilde{\pi})^{-1}\bigl(\hat{t}\bigr)\cong \operatorname{CP}^2 \#9\overline{\operatorname{CP}}^2$} is a compact complex manifold for $\hat{t} \in T$, $(\mathcal{F},\mathcal{S},T,\widetilde{\pi},o,j)$ is a smooth compactifiable deformation.

Moreover, the following lemma from Ehresmann indicates that $(\mathcal{F},\mathcal{S},T,\widetilde{\pi},o,j)$ is differentially trivial along $T$.

\begin{Lemma}[{\cite[Theorem~9.3]{Claire}}] \label{Lemma3.15}
 For any proper holomorphic submersion $g\colon M\rightarrow N$ with $M$ and $N$ being complex manifolds, $g$ is a locally trivial fibration.
\end{Lemma}

Taking use of Lemma~\ref{Lemma3.15}, since $\widetilde{\pi}|_{\mathcal{F}\backslash \mathcal{S}}$ is a proper submersion from the complex manifold~$\mathcal{F}\backslash \mathcal{S}$ to the complex manifold $T$, $\widetilde{\pi}|_{\mathcal{F}\backslash \mathcal{S}}$ is a locally trivial fibration. So $(\mathcal{F},\mathcal{S},T,\widetilde{\pi},o,j)$ is differentially trivial along $T$.

In conclusion, $(\mathcal{F},\mathcal{S},T,\widetilde{\pi},o,j)$ is a smooth compactifiable deformation differentially trivial along $T$. So Proposition \ref{Proposition3.14} is proved.
\end{proof}


Therefore, $(\mathcal{F},\mathcal{S},T,\widetilde{\pi},o,j)$ is a compactifiable deformation of the quasi-projective variety $S\backslash C$ over a 9-dimensional complex manifold $T$ with $o=(\tau_0,q_1,q_2,\dots ,q_8)$ and $j\colon S \backslash C \rightarrow (\widetilde{\pi}^{-1}(o))$ defined by $j(x)=(x,\tau_0,q_1,q_2,\dots ,q_8)$ for $x \in S\backslash C$. Combining the results of Theorem~\ref{Theorem3.13} and Proposition \ref{Proposition3.14}, the proof of Theorem~\ref{Theorem1.1} is completed.

\subsection[Complete K\"ahler metrics on each fiber of the deformation family of quasi-projective varieties]{Complete K\"ahler metrics on each fiber of the deformation family\\ of quasi-projective varieties}

The following is to construct the complete K\"ahler metrics on each fiber of the deformation family of quasi-projective varieties. A key point for the construction here is to take use of Macro Brunella's statement in \cite{Marco}.

Taking use of the notations from the introduction and Section~\ref{sec2}, the complete K\"ahler metrics will be constructed on $S(\tau)\backslash C(\tau)$ for $\tau \in Y$.

For any integers $k\geq 2$ and $d \geq 3k+1$ satisfying $10 \leq d^2/k^2$, let $L=d\cdot H-k\cdot\sum _{i=1}^9 E_i$. Taking use of Lemma~\ref{Lemma2.5}, $L$ is an ample line bundle on $S(\tau)$. Then there exists $n\in \mathbb{N}$ such that~$L^n$ and $L^n\otimes [-C(\tau)]$ are very ample \cite{Phillip}. Now let $ \{ g_0,g_1,\dots ,g_{\theta} \}$ be a basis for the space of holomorphic sections on $L^n\otimes [-C(\tau)]$ with $\theta \in \mathbb{N}$. Let $\hat{s}$ be the defining section of $C(\tau)$ and~$ \{ \zeta \}$ be a frame for $L$. Then $ \{ g_0\hat{s},g_1\hat{s},\dots ,g_{\theta} \hat{s} \}$ are all holomorphic sections on $L^n$ with value zero on~$C(\tau)$. This is an important point for the construction. Then the hermitian metric $h_L$ defined on $L|_{S(\tau)\backslash C(\tau)}$ through the $C^\infty $ function
\[
h_L(x)=\frac{\zeta(x) \cdot \overline{\zeta (x)}}{\bigl(|g_0(x)\cdot \hat{s}(x)|^2+|g_1(x)\cdot \hat{s}(x)|^2+\dots +|g_{\theta}(x)\cdot \hat{s}(x)|^2\bigr)^{\frac{1}{n}}}
\]
for $x \in S(\tau)\backslash C(\tau)$ is smooth.

\begin{Lemma} \label{Lemma3.16}
The curvature form $\Theta_{h_L}$ of $h_L$ is positive.
\end{Lemma}

\begin{proof}
Since $L^n \otimes [-C(\tau)]$ is very ample, taking use of the Kodaira embedding theorem \cite{Phillip}, the evaluation map
\[\Phi_{L^n \otimes [-C(\tau)]}(x)=[g_0(x):g_1(x):\dots :g_{\theta}(x)]\in \operatorname{CP}^{\theta}\]
gives a holomorphic embedding of \smash{$S(\tau) \cong \operatorname{CP}^2 \#9\overline{\operatorname{CP}}^2$} into $\operatorname{CP}^{\theta}$. In addition, let $[\xi_0,\xi_1,\dots ,\xi_{\theta}]$ be the coordinate system of $\operatorname{CP}^{\theta}$ coinciding with the image of the evaluation map $\Phi_{L^n \otimes [-C(\tau)]}$.

The functions
\[
g_i([\xi_0:\xi_1:\dots :\xi_{\theta}])=\frac{|\xi_i|^2}{|\xi_0|^2+|\xi_1|^2+\dots +|\xi_{\theta}|^2}
\]
on $\big \{ [\xi_0:\xi_1:\dots :\xi_{\theta}]\in \operatorname{CP}^\theta\mid \xi_i\neq 0\big \}$ for $i \in \{ 0,1,\dots ,\theta \}$ define a hermitian metric $g$ on the hyperplane line bundle $\mathcal{O}(1)$. Then the curvature form for the hermitian metric $g$ is
\[
\Theta_g=\frac{\sqrt{-1}}{\pi}\partial\overline{\partial}\log\bigl(|\xi_0|^2+|\xi_1|^2+\dots +|\xi_{\theta}|^2\bigr)
\]
for $[\xi_0:\xi_1:\dots :\xi_{\theta}]\in \operatorname{CP}^\theta$. Furthermore, the associate $(1,1)$-form for the Fubini--Study metric on $\operatorname{CP}^\theta$ is the positive K\"ahler form
\[\frac{\sqrt{-1}}{2\pi}\partial\overline{\partial}\log\bigl(|\xi_0|^2+|\xi_1|^2+\dots +|\xi_{\theta}|^2\bigr)\]
for $[\xi_0:\xi_1:\dots :\xi_{\theta}]\in \operatorname{CP}^\theta$. So $\Theta_g$ is a positive form on $\operatorname{CP}^\theta$.

Therefore, $g$ is a smooth positive hermitian metric on the hyperplane line bundle $\mathcal{O}(1)$ for the coordinate system $[\xi_0:\xi_1:\dots :\xi_{\theta}]$ on $\operatorname{CP}^{\theta}$. Then the pullback metric $\hat{g}=(\Phi_{L^n \otimes [-C(\tau)]})^*g$ is a smooth positive hermitian metric on the pullback bundle $(\Phi_{L^n \otimes [-C(\tau)]})^*\mathcal{O}(1)=L^n \otimes [-C(\tau)]$. So the curvature form $\Theta_{\hat{g}}$ of $\hat{g}$ is positive.

Moreover, in \cite{Marco}, Macro Brunella stated that under the condition that the normal bundle~\smash{$N_{C(\tau)/S(\tau)}$} satisfied the Diophantine condition, \smash{$K^{-1}_{S(\tau)}=[C(\tau)]$} admitted a smooth hermitian metric with semi-positive curvature. So $\log|\hat{s}|$ gives the local weights of a (singular) metric $g_{\hat{s}}$ on $K^{-1}_{S(\tau)}$ with $\Theta_{g_{\hat{s}}}>0$ \cite[Definition~8.7]{Vincent-Ahmed}. Therefore, $\Theta_{h_L}=\frac{1}{n}(\Theta_{\hat{g}}|_{S(\tau)\backslash C(\tau)}+\Theta_{g_{\hat{s}}}|_{S(\tau)\backslash C(\tau)})$ is positive. Then Lemma~\ref{Lemma3.16} is proved. \end{proof}


In addition, taking use of Theorem~\ref{Theorem2.19}, let $L|_{C(\tau)}$ be the restricted line bundle of $L$ on $C(\tau)$. Furthermore, let $L^*$ be the pull back bundle of $L|_{C(\tau)}$ on $V_{0,\infty}$. Then $L^*|_{V_{0,r}}=L|_{V_{0,r}}$. Taking use of the equivalent condition in \cite[Theorem]{Christian} (also see \cite{Takayuki-Takato2}), there exists a corresponding theta line bundle on $U_0=\mathbb{C}^2/{\Lambda_0}$ holomorphically isomorphic to the pull back line bundle $(f_{U_0})^*(L^*)$.

Let \smash{$H_1=\bigl(\begin{smallmatrix}
\frac{(L.C(\tau))}{\operatorname{ Im}\tau} & 0\\
0&0
\end{smallmatrix}\bigr)$}. Here, $H_1(x,y)=x^tH_1 \overline{y}$ for $x,y \in \mathbb{C}^2$. Defining
$\rho\colon \Lambda_0 \rightarrow U(1)$ through
\[
\rho(\lambda+\mu)=\rho(\lambda)\rho(\mu){\rm e}^{\pi\sqrt{-1}\cdot\operatorname{ Im}H_1(\lambda,\mu)}
\]
 for $ \lambda, \mu \in \Lambda_0 $, let \[\alpha_\lambda(x)=\rho (\lambda){\rm e}^{\pi H_1(x,\lambda)+(\pi/2)\cdot H_1(\lambda,\lambda)}\] for $ \lambda \in \Lambda_0 $ and $ x \in \mathbb{C}^2 $.

The theta line bundle on $U_0=\mathbb{C}^2/{\Lambda_0}$ corresponding to \smash{$H_1=\bigl(\begin{smallmatrix}
\frac{(L.C(\tau))}{\operatorname{ Im} \tau} & 0\\
0&0
\end{smallmatrix}\bigr)$} can be defined~as ${
L_{H_1,\rho}=\bigl(\mathbb{C}_{\zeta} \times \mathbb{C}^2\bigr)/{\Lambda_0}}
$
with $\lambda \cdot (\zeta, x)=$$(\alpha_\lambda(x)\cdot \zeta, x+\lambda )$ for $\lambda \in\Lambda_0$, $\zeta \in \mathbb{C}_\zeta$ and $x \in \mathbb{C}^2$.

\begin{Lemma} \label{Lemma3.17}
The theta line bundle $L_{H_1,\rho}$ on $U_0=\mathbb{C}^2/{\Lambda_0}$ is holomorphically isomorphic to the pull back bundle $(f_{U_0})^*(L^*)$ with \smash{$H_1=\bigl(\begin{smallmatrix}
\frac{(L.C(\tau))}{\operatorname{ Im} \tau} & 0\\
0&0
\end{smallmatrix}\bigr)$}.
\end{Lemma}

\begin{proof}
For the pull back bundle \smash{$L^*_{H_1,\rho}=\bigl(f^{-1}_{U_0}\bigr)^*L_{H_1,\rho}$} on $V_{0,\infty}$, according to \cite[Proposition~3.5]{Takayuki-Takato2}, $\exists \hat{G} \in \operatorname{Pic}(W) $ such that \smash{$\hat{G}|_{V_{0,r}}=L^*_{H_{1,\rho}}|_{V_{0,r}}$}. Then \smash{$(\hat{G}.C(\tau))=(L.C(\tau))$} (see \cite[Lemma~3.4]{Takayuki-Takato2}).

Moreover, according to \cite[Lemma~3.4]{Takayuki-Takato2} and \cite[Proposition~3.5]{Takayuki-Takato2}, $L_{H_1,\rho}$ with $H_1=\smash{\bigl(\begin{smallmatrix}
\frac{(L.C(\tau))}{\operatorname{ Im}\tau} & 0\\
0&0
\end{smallmatrix}\bigr)}$ is the unique theta line bundle defined in the formula as above such that $\hat{G}$ has the same intersection number with $C(\tau)$ as $L$. Therefore, the theta line bundle $L_{H_1, \rho}$ on $U_0=\mathbb{C}^2_{(z,\eta)}/{\Lambda_0}$ is holomorphically isomorphic to the pull back bundle $(f_{U_0})^*(L^*)$ with \smash{$H_1\!=\!\bigl(\!\begin{smallmatrix}
\frac{(L.C(\tau))}{\operatorname{ Im}\tau} & 0\\
0&0
\end{smallmatrix}\bigr)$}. So Lemma~\ref{Lemma3.17} is proved.
\end{proof}


Taking use of Lemma~\ref{Lemma3.17},
\[
c_1(L|_{C(\tau)})=\left[\frac{(L.C(\tau))}{\operatorname{Im}\tau}\cdot \sqrt{-1}\text{{\rm d}z}\land {\rm d}\bar{z}\right].
\]
In addition, according to \cite[Lemma~7.31]{Vincent-Ahmed} ($\partial \overline{\partial}$-lemma) and \cite[Section 8.2]{Vincent-Ahmed}, there exists a~smooth hermitian metric $h_{C(\tau)}$ on $L|_{C(\tau)}$ such that
\[
\Theta_{h_{C(\tau)}} = \frac{(L.C(\tau))}{\operatorname{ Im}\tau} \sqrt{-1} {\rm d}z \wedge {\rm d} \bar{z}.
\]
 So $(\pi_W)^*h_{C(\tau)}$ is a smooth hermitian metric on $L|_W$ such that
\[
\Theta_{(\pi_W)^*h_{C(\tau)}}=\frac{(L.C(\tau))}{\operatorname{Im}\tau} \sqrt{-1} {\rm d}z \wedge d \bar{z}.
\]

For the line bundle \smash{$L|_{C(\tau)}$}, there exist open finite covers \smash{$\{ \mathcal{U}_{\alpha_1} \}_{\alpha_1 \in I}$} of $C(\tau)$ and local trivializations
\[
\Phi_{\alpha_1}\colon \ \pi^{-1}_{L|_{C(\tau)}}(\mathcal{U}_{\alpha_1}) \rightarrow \mathcal{U}_{\alpha_1} \times \mathbb{C}
\]
 mapping \smash{$\pi^{-1}_{L|_{C(\tau)}}(x)$} isomorphically onto $ \{ x \} \times \mathbb{C}$ for $x \in \mathcal{U}_{\alpha_1}$ and $\alpha_1 \in I$ with $I$ being an index set. Let
\[
g_{\alpha_1 \beta_1}(x)=\big[\Phi_{\alpha_1} \circ (\Phi_{\beta_1})^{-1}\big]|_{ \{ x \} \times \mathbb{C}}
\]
for $ x \in \mathcal{U}_{\alpha_1} \cap \mathcal{U}_{\beta_1}$ with $\alpha_1, \beta_1 \in I$.

The line bundle \smash{$L|_W=(\pi_W)^*L|_{C(\tau)}$} and \smash{$\big \{ W_{\alpha_1}=\pi^{-1}_W(\mathcal{U}_{\alpha_1})\big \}_{\alpha_1 \in I}$} is an open cover of $W$. Moreover, there exist local trivializations
\[
\Psi_{\alpha_1}\colon \ \pi^{-1}_{L|_W}(W_{\alpha_1})\rightarrow W_{\alpha_1} \times \mathbb{C}
\]
 mapping \smash{$\pi^{-1}_{L|_W}(x)$} isomorphically onto $\{ x \} \times \mathbb{C}$ induced by $\Phi_{\alpha_1}$ for $x=[(z,w)] \in W_{\alpha_1}$ and $\alpha_1 \in I$. Let
\[
\widetilde{g}_{\alpha_1 \beta_1}(x)=\big[\Psi_{\alpha_1} \circ (\Psi_{\beta_1})^{-1}\big]\big|_{\{ x \}\times \mathbb{C}}
\]
for $x \in W_{\alpha_1} \cap W_{\beta_1}$ with $\alpha_1,\beta_1 \in I$.

A singular hermitian metric for a line bundle is defined through local weights \cite[Definition~8.7]{Vincent-Ahmed}. For $\varepsilon > 0$, let
$
 \varphi^{\alpha_1}_L \in \operatorname{Psh}(W_{\alpha_1} \backslash \mathcal{U}_{\alpha_1})\cap C^\infty(W_{\alpha_1} \backslash \mathcal{U}_{\alpha_1})$ and $\varphi^{\alpha_1} _{C(\tau)}+\log \varepsilon \in \operatorname{Psh}(W_{\alpha_1}) \cap C^\infty (W_{\alpha_1})
$
be the local weights of the metrics $h_L$ and \smash{$\varepsilon^{-1}(\pi_W)^*h_{C(\tau)}$} with $\operatorname{Psh}(W_{\alpha_1} \backslash \mathcal{U}_{\alpha_1})$ and $\operatorname{Psh}(W_{\alpha_1})$ being the sets consisting of plurisubharmonic functions on $W_{\alpha_1} \backslash \mathcal{U}_{\alpha_1}$ and $W_{\alpha_1}$ for $\alpha_1 \in I$, respectively.

Here, it is another important point for the construction modified from \cite{Takayuki-Takato2}. Through the following theorem, $\varphi^{\alpha_1}_L$ and $\varphi^{\alpha_1}_{C(\tau)}+\log \varepsilon$ can be patched to smooth plurisubharmonic functions on $W_{\alpha_1}$ for $\alpha_1 \in I$.

Let $\gamma =(\gamma_1,\gamma_2,\dots ,\gamma_p)$ with $\gamma_i >0$ for $i \in \left \{ 1,\dots ,p\right \}$ and $X$ be a compact complex manifold. Moreover, let $\hat{\theta}$ be a real non-negative smooth function defined on $ \mathbb{R}$ supported on $[-1,1]$ with
\[
\int _\mathbb{R} \hat{\theta}(x){\rm d}x=1 \qquad \text{and} \qquad \int _\mathbb{R}x\hat{\theta} (x){\rm d}x=0.
\]
 In addition, let
\[
M_\gamma (t_1,t_2,\dots ,t_p)=\int _{\mathbb{R}^p}\max \{ t_1+h_1,\dots ,t_p+h_p\}\prod_{1 \leq j \leq p} \gamma^{-1}_j \hat{\theta} (h_j/\gamma_j) {\rm d}h_1 {\rm d}h_2 \dots .{\rm d} h_p
\]
for $(t_1,t_2,\dots ,t_p) \in \mathbb{R}^p$ which is called a regularized max function \cite{Alexandre}.

\begin{Lemma}[{\cite[Lemma~5.18]{Jean}}] \label{Lemma3.18}
 $M_\gamma (t_1,t_2,\dots ,t_p)$ defined above satisfying the following properties:
\begin{itemize}\itemsep=0pt
\item[$(1)$]$M_\gamma (t_1,t_2,\dots ,t_p)$ is smooth and convex in $\mathbb{R}^p$;
\item[$(2)$]$M_\gamma (t_1+a,t_2+a,\dots ,t_p+a)=M_\gamma (t_1,t_2,\dots ,t_p)+a$, for any $a \in \mathbb{R}$ and $(t_1,t_2,\dots ,t_p) \in \mathbb{R}^p$.
\end{itemize}
\end{Lemma}

Let $h=h_L$ on $S(\tau)\backslash W$. Moreover, let
\[
M_{(1,1)}\bigl(\varphi^{\alpha_1}_L,\varphi^{\alpha_1}_{C(\tau)}+\log\varepsilon\bigr)=\int_{\mathbb{R}^2}\max\bigl\{\varphi^{\alpha_1}_L+h_1,\varphi^{\alpha_1}_{C(\tau)}
+\log\varepsilon+h_2 \bigr\}\prod_{1\leq j\leq 2}\hat{\theta}(h_j){\rm d}h_1{\rm d}h_2\]
be the local weight of the metric $h$ on $W_{\alpha_1}$ for $\alpha_1 \in I$.

Let $r_1$ and $r_2$ be positive numbers with $r_1 <r_2<r$. Choosing $\varepsilon$ small enough,
\begin{gather*}
M_{(1,1)}\bigl(\varphi^{\alpha_1}_L,\varphi^{\alpha_1}_{C(\tau)}+\log \varepsilon\bigr)\\
\qquad=\int_{\mathbb{R}^2}\max \bigl\{\varphi^{\alpha_1}_L+h_1,\varphi^\alpha_{C(\tau)}+\log\varepsilon +h_2 \bigr\} \prod_{1 \leq j \leq 2} \hat{\theta} (h_j){\rm d} h_1 {\rm d} h_2 =\varphi^{\alpha_1}_L
\end{gather*}
on $W_{\alpha_1} \cap \{ [(z,w)] \in W \mid r_1 <|w| \}$ for $\alpha_1 \in I$. That is to say, $h=h_L$ on $ \{ [(z,w)] \in W \mid r_1<|w| \}$.

Taking use of property~(1) of Lemma~\ref{Lemma3.18}, $M_{(1,1)}\bigl(\varphi^\alpha_L,\varphi^\alpha_{C(\tau)}+\log \varepsilon\bigr)$ is plurisubharmonic for~${\alpha\in I}$. Furthermore, from property~(2) of Lemma~\ref{Lemma3.18}, $h$ is a smooth hermitian metric on~$L$. Therefore, $h$ is a smooth hermitian metric on~$L$ with $\Theta_{h}\geq 0 $ (see \cite[Section~8.2]{Vincent-Ahmed} for the definition of singular hermitian metrics).

Furthermore, there exists $\varepsilon_0 >0$ satisfying $\sqrt{\varepsilon_0}r<r_1$ such that
\[
M_{(1,1)}\bigl(\varphi^{\alpha_1}_L,\varphi^{\alpha_1}_{C(\tau)}+\log \varepsilon\bigr)=\varphi^{\alpha_1}_L
\]
on $W_{\alpha_1} \cap \{ [(z,w)] \in W\mid |w|< \sqrt{\varepsilon_0}r \}$ for $\alpha_1 \in I$. So
$h=\varepsilon^{-1}(\pi_W)^*h_{C(\tau)}$ on $\{ [(z,w)] \in W \mid |w|< \sqrt{\varepsilon_0} r \}$.

Let $\theta_s(t)=\bigl(\log\bigl(t^2/|s|\bigr)\bigr)^2$ be a smooth function for $t>0$ and $s \in \{ s\mid s \in \Delta \backslash \{ 0 \},\, |s|<\varepsilon_0 \}$ with $\Delta = \{ s \in \mathbb{C} \mid |s|<1 \}$.

Let $f(x)=1$ for $|x|<r-\frac{r-r_2}{2}$ and $f(x)=0$ for $r-\frac{r-r_2}{2} \leq|x|$. The standard mollifier~${\eta \in C^\infty(\mathbb{R})}$ is defined as
\[\eta(x):=\begin{cases}
a\cdot {\rm e}^{1/(|x|^2-1)} & \text{if}\  |x|<1, \\
0 & \text{if}\  |x| \geq 1,
\end{cases}\]
with \smash{$a=1/\int^{1}_{-1}{\rm e}^{1/(|x|^2-1)}{\rm d}x>0$} \cite{Lawrence}.

Moreover, let $\eta_\delta(x):=\frac{1}{\delta}\eta(\frac{x}{\delta})$. Then for $\delta=\frac{r-r_2}{4}$, let
\[\widetilde{f}(x)=\eta_\delta * f(x)=\int^\infty_{-\infty}\eta_\delta(x-y)f(y){\rm d}y.\]
Then $\widetilde{f}$ is a smooth function on $\mathbb{R}$ with $\widetilde{f}(x)=1$ for $|x|<r_2$ and $\widetilde{f}(x)=0$ for $r<|x|$.

Let $\Psi_s\colon S(\tau) \backslash C(\tau) \rightarrow \mathbb{R}$ be the smooth function defined by
\[\Psi_s(p):= \begin{cases}
\widetilde{\theta}_s(|w|), & p=[(z,w)] \in W\backslash C(\tau), \\
\widetilde{\theta}_s(r), & p \notin W
\end{cases}\]
with $\widetilde{\theta}_s(x)=\widetilde{f}(x) \cdot \theta_s(x)$ for $x>0$.

Then the following is to give the modified complete K\"ahler metrics on $S(\tau)\backslash C(\tau)$ analogous to the complete K\"ahler metrics described in \cite{Takayuki-Takato2}.

\begin{Theorem}[{\cite{Takayuki-Takato2}}] \label{Theorem3.19}
 There exists $b>0$ such that the metric $h \cdot {\rm e}^{-b \Psi_s}$ is a smooth hermitian metric on $L|_{S(\tau)\backslash C(\tau)}$ with $\Theta_{h \cdot {\rm e}^{-b \Psi_s}}>0$. Moreover, the K\"ahler form $\Theta_{h \cdot {\rm e}^{-b \Psi_s}}$ is of the form
\[\omega|_{W_{\varepsilon_0}}=\frac{ (L.C(\tau))}{\operatorname{ Im}\tau}\cdot \sqrt{-1} {\rm d}z \wedge {\rm d} \bar{z}+\frac{2b}{\pi} \cdot \frac{\sqrt{-1}{\rm d}w \wedge {\rm d} \bar{w}}{|w|^2}\]
on the set $W_{\varepsilon_0}=\{ [(z,w)]\in W \backslash C(\tau)\mid |w|<\sqrt{\varepsilon_0}r \}$. In addition, $\omega|_{W_{\varepsilon_0}}$ is Ricci-flat. Then the K\"ahler form $\Theta_{h\cdot {\rm e}^{-b\Psi_s}}$ gives a complete K\"ahler metric on $S(\tau)\backslash C(\tau)$.
\end{Theorem}

The complete proof of Theorem~\ref{Theorem3.19} can be found in Appendix~\ref{appendixA}. Therefore, the quasi-projective variety above tends to be a non-compact complete K\"ahler manifold.

\subsection{A deformation family of non-compact complete K\"ahler manifolds}

The following is to show that the deformation family $(\mathcal{F},\mathcal{S},T,\widetilde{\pi},o,j)$ can be easily changed to a deformation family satisfying the conditions in Definition~\ref{Definition3.5}.

Since the fibers are now proved to be non-compact complete K\"ahler manifolds, the deformation family above is proved to be a deformation family of non-compact complete K\"ahler manifolds over a 9-dimensional complex manifold.
Moreover, for $T$ defined above, there exists a complex disc $D_3 \subset T$ centred at $o=(\tau_0,q_1,q_2,q_3,q_4,q_5,q_6,q_7,q_8)$.

\begin{Corollary} \label{Corollary3.20}
Let $f\colon T \rightarrow \mathbb{C}^9$ be defined by $f(x)=x-o$. Then $f \circ \widetilde{\pi}|_{\mathcal{F}\backslash \mathcal{S}}\colon\mathcal{F} \rightarrow f(D_3)$ is a~deformation family of the non-compact complete K\"ahler manifold $j(S\backslash C)$ over a 9-dimensional complex manifold satisfying the conditions in Definition {\rm\ref{Definition3.5}}.
\end{Corollary}

The complete proof of Corollary~\ref{Corollary3.20} can be found in Appendix~\ref{appendixA}. Then the construction of the deformation family of the non-compact complete K\"ahler manifold is completed.

\section{A deformation family of symmetric projective K3 surfaces}\label{sec4}

\subsection{Symmetric K\"ahler metrics on the symmetric projective K3 surfaces}

The following is to give symmetric K\"ahler metrics on the symmetric projective K3 surfaces constructed in Section~\ref{sec2}.

In Section~\ref{sec3}, the complete K\"ahler metric is constructed on $S(\tau)\backslash C(\tau)$ for $\tau \in Y$. That is to say, the complete K\"ahler metrics are constructed on fibers of the smooth compactifiable deformation of $S\backslash C$.

Taking use of the notations from Section~\ref{sec2}, the symmetric projective K3 surface is constructed in Theorem~\ref{Theorem2.8} through $S^+(\tau)\backslash C^+(\tau)$ and $S^-(\tau)\backslash C^-(\tau)$.

Since $S^+(\tau)\backslash C^+(\tau)$ is just a copy of $S(\tau)\backslash C(\tau)$, let $\omega^+$ be a copy of $\omega$ on $S(\tau)\backslash C(\tau)$ constructed above. Let $L^{+}=d\cdot H^{+}-k \cdot \sum _{i=1}^9 E_i^{+}$ be defined on $S^+(\tau)$ as a copy of $L$ in Section~\ref{sec3}. Then \smash{$L^{-}=d\cdot H^{-}-k\cdot\sum _{i=1}^9 E_i^{-}$} is an ample line bundle on $S^-(\tau)$ and \smash{$(\ell_{S^+(\tau)})^*L^-=L^+$}. In addition, let $h^+$ be defined on $L^+$ as a copy of $h\cdot {\rm e}^{-b\Psi_s}$ on $L$. Then there is a hermitian metric~$h^-$ defined on $L^-$ satisfying \smash{$h^+=(\ell_{S^+(\tau)\backslash C^+(\tau)})^*h^-$} such that \smash{$\omega^-=\Theta_{h^-}=\bigl(\ell^{-1}_{S^+(\tau)\backslash C^+(\tau)}\bigr)^* \omega^+$} is a~K\"ahler form on $S^-(\tau)\backslash C^-(\tau)$.

Identifying $V_s^+$ and $V_s^-$ through $f_s$, $M^-_s$ and $M^+_s$ can be glued to a K3 surface $X_s$ together with a line bundle $L_s$ derived from the line bundles $L^+$ and $L^-$. Taking use of the result from Section~\ref{sec3}, $\exists \ \varepsilon_0 >0$ such that $\omega$ is well defined on $S(\tau)\backslash C(\tau)$ with the results of Theorem~\ref{Theorem3.19} holding for $s \in \{ s\mid s \in \Delta \backslash \{ 0 \},\, |s|<\varepsilon_0 \}$. Then $\omega^+$ and $\omega^-$ can be glued to be a K\"ahler form $\omega_s$ on $X_s$ for $s \in \{ s\mid s \in \Delta \backslash \{ 0 \},\, |s|<\varepsilon_0 \}$. That is to say, $L_s$ is ample for~${s \in \{ s\mid s \in \Delta \backslash \{ 0 \},\, |s|<\varepsilon_0 \}}$. So $X_s$ is projective for $s \in \{ s\mid s \in \Delta \backslash \{ 0 \}, \, |s|<\varepsilon_0 \}$ (also see~\cite[Theorem~1.2]{Takayuki-Takato2}).

\begin{Definition} \label{Definition4.1}
Let $\mathcal{K}$ be a K3 surface with an ample line bundle $L_\mathcal{K}$ and a non-trivial holomorphic involution $f_\mathcal{K}$ satisfying $f^*_{\mathcal{K}}L_{\mathcal{K}}=L_{\mathcal{K}}$. Then if $\widehat{\omega}$ is a K\"ahler metric on $\mathcal{K}$ such that~${f^*_{\mathcal{K}}\widehat{\omega}=\widehat{\omega}}$, then $\widehat{\omega}$ is said to be a symmetric K\"ahler form on $\mathcal{K}$ corresponding to a symmetric K\"ahler metric.
\end{Definition}

With the discussion above, the following lemma is true.

\begin{Lemma} \label{Lemma4.2}
There exists a symmetric K\"ahler metric on the projective symmetric K3 surface~$X_s$ for $s \in \{ s\mid s \in \Delta \backslash \{ 0 \},\, |s|<\varepsilon_0 \}$.
\end{Lemma}

\begin{proof}
For $s \in \{ s\mid s \in \Delta \backslash \{ 0 \},\, |s|<\varepsilon_0 \}$, since $\omega^+=(\ell_{S^+(\tau)\backslash C^+(\tau)})^*\omega^-$, $\omega_s=F^*_{X_s}\omega_s$. Then~$\omega_s$ is the symmetric K\"ahler form needed for $s \in \{ s\mid s \in \Delta \backslash \{ 0 \},\, |s|<\varepsilon_0 \}$. So there exists a~symmetric corresponding K\"ahler metric on the projective symmetric K3 surface $X_s$ for $s \in \{ s\mid s \in \Delta \backslash \{ 0 \},\, |s|<\varepsilon_0 \}$. Then Lemma~\ref{Lemma4.2} is proved.
\end{proof}

\subsection{A deformation family of symmetric projective K3 surfaces}

The following is to give a deformation family of symmetric projective K3 surfaces over a 10-dimensional complex manifold and complete the proof of Corollary~\ref{Corollary1.2}. With the ample line bundles provided above, the proofs can be completed directly through the maps derived.

Let $\hat{\Delta}=\{ s\mid s \in \Delta \backslash \{ 0 \},\, |s|<\varepsilon_0 \}$. Then there are manifolds
\[\mathcal{M}^+=\bigl(S(\tau)\times \hat{\Delta}\bigr)\backslash \bigl\{ ([(z,w)],s)\in W \times \hat{\Delta} \mid |w|\leq \sqrt{|s|}/r \bigr\}=\bigl\{ (x,s) \in M^+_s \times \hat{\Delta}\bigr\}\]
and
$\mathcal{M}^-=\bigl\{(\ell_{S^+(\tau)}(x),s) \mid (x,s) \in M^+_s \times \hat{\Delta}\bigr\}$.
In addition, let
\[\mathcal{V}^+=\bigl\{ ([(z,w)],s)\in W \times \hat{\Delta} \mid \sqrt{|s|}/r <|w| <\sqrt{|s|} r\bigr\}=\bigl\{ (x,s) \in V^+_s \times \hat{\Delta}\bigr\}\]
and
$\mathcal{V}^-=\bigl\{ \bigl(\ell_{S^+(\tau)}(x),s\bigr) \mid (x,s) \in V^+_s \times \hat{\Delta}\bigr\}$
be the submanifolds of $\mathcal{M}^+$ and $\mathcal{M}^-$, respectively.

Define $f\colon\mathcal{V}^+ \rightarrow \mathcal{V}^-$ by
$f([(z,w)],s)=([(g_\xi(z),s/w)],s)=(f_s([(z,w)]),s)$
for $([(z,w)],s)\allowbreak\in \mathcal{V}^+$. Then through the map $f$, $\mathcal{M}^+$ and $\mathcal{M}^-$ can be glued to a complex manifold~$\mathcal{M}$. Let~${\pi_{\mathcal{M}}\colon\mathcal{M} \rightarrow \hat{\Delta}}$ be the natural projection. Then $\bigl(\mathcal{M}, \hat{\Delta}, \pi_{\mathcal{M}}\bigr)$ is a complex analytic family of symmetric projective K3 surfaces over the one dimensional complex manifold~\smash{$\hat{\Delta}$} (also see \cite{Takayuki-Takato2} for a more general result).

For each fiber of the smooth compactifiable deformation $(\mathcal{F},\mathcal{S},T,\widetilde{\pi},o,j)$ constructed in Theorem~\ref{Theorem1.1}, a complex analytic family of symmetric projective K3 surfaces over one dimensional complex manifold $\hat{\Delta}=\{ s\mid s \in \Delta \backslash \{ 0 \},\, |s|<\varepsilon_0 \}$ can be constructed.

In fact, there is an analogous deformation family of quasi-projective varieties constructed here. Let $\mathcal{F}^-$ be the result manifold after blowing up $\operatorname{CP}^2\times T$ at the points
\begin{gather*}
\bigcup_{(\widetilde{\tau},\vec{x})=(\widetilde{\tau},x_1,\dots ,x_8)\in T}\{(\ell_{CP^2}\circ f_{\widetilde{\tau}}\circ f_{\widetilde{\tau},1}(x_1),\widetilde{\tau},\vec{x}),\dots ,\\
\hphantom{\bigcup_{(\widetilde{\tau},\vec{x})=(\widetilde{\tau},x_1,\dots ,x_8)\in T}} (\ell_{CP^2}\circ f_{\widetilde{\tau}}\circ f_{\widetilde{\tau},8}(x_8),\widetilde{\tau},\vec{x}),(\ell_{CP^2}\circ f_{\widetilde{\tau}}\circ f_9(\widetilde{\tau},\vec{x}),\widetilde{\tau},\vec{x}\}.
\end{gather*}
Let $\mathcal{S}^-$ be the analogous complex manifold of $\mathcal{S}$ and $\widetilde{\pi}^-\colon\mathcal{F}^- \rightarrow T^-=T$ be the map induced by $\widetilde{\pi}$. Then $(\mathcal{F}^-, T^-, \widetilde{\pi}^-)$ is a deformation family of \smash{$\operatorname{CP}^2 \#9\overline{\operatorname{CP}}^2$} analogous to the complex analytic family $(\mathcal{F}, T, \widetilde{\pi})$. Moreover, $\widetilde{\pi}^-|_{\mathcal{F}^-\backslash \mathcal{S}^-}\colon{\mathcal{F}^-\backslash \mathcal{S}^-} \rightarrow T^-$ gives a deformation family of quasi-projective varieties analogous to $(\mathcal{F},\mathcal{S},T,\widetilde{\pi},o,j)$.

In addition, let $\bigl(\mathcal{F}^+, T^+, \widetilde{\pi}^+\bigr)$ be a copy of the deformation family $(\mathcal{F}, T, \widetilde{\pi})$ and $\mathcal{S}^+$ be a copy of $\mathcal{S}$. Then the following theorem is true.

\begin{Theorem} \label{Theorem4.3}
There is a deformation family of symmetric projective K3 surfaces over a $10$-dimensional complex manifold constructed from $(\mathcal{F},\mathcal{S},T,\widetilde{\pi},o,j)$ and an analogous deformation family.
\end{Theorem}

\begin{proof}
Here, $(\mathcal{F},\mathcal{S},T,\widetilde{\pi},o,j)$ is a smooth compactifiable deformation of $S \backslash C$ over a 9-di\-men\-sion\-al complex manifold. $\widetilde{\pi}^-|_{\mathcal{F}^-\backslash \mathcal{S}^-}\colon {\mathcal{F}^-\backslash \mathcal{S}^-} \rightarrow T^-$ gives a deformation family of quasi-projective varieties analogous to $(\mathcal{F},\mathcal{S},T,\widetilde{\pi},o,j)$ and $\widetilde{\pi}^+|_{\mathcal{F}^+\backslash \mathcal{S}^+}\colon{\mathcal{F}^+\backslash \mathcal{S}^+} \rightarrow T^+$ gives a deformation family as a copy of $(\mathcal{F},\mathcal{S},T,\widetilde{\pi},o,j)$.

From the discussion above, choosing two corresponding fibers from $\mathcal{F}^+$ and $\mathcal{F}^-$, respectively, a~deformation family of symmetric projective K3 surfaces over one dimensional complex manifold~$\hat{\Delta}=\{ s\mid s \in \Delta \backslash \{ 0 \}, \, |s|<\varepsilon_0 \}$ can be constructed. Therefore, a deformation family of symmetric projective K3 surfaces over a 10-dimensional complex manifold can be constructed through the deformation families induced by $\widetilde{\pi}^-|_{\mathcal{F}^-\backslash \mathcal{S}^-}\colon{\mathcal{F}^-\backslash \mathcal{S}^-} \rightarrow T^-$ and~$\widetilde{\pi}^+|_{\mathcal{F}^+\backslash \mathcal{S}^+}\colon{\mathcal{F}^+\backslash \mathcal{S}^+} \rightarrow T^+$. So Theorem~\ref{Theorem4.3} is proved.
\end{proof}

Therefore, combining the results from Lemma~\ref{Lemma4.2} and Theorem~\ref{Theorem4.3}, the proof of Corollary~\ref{Corollary1.2} is completed.

\appendix

\section{Appendix}\label{appendixA}

The main goal of this appendix is to give the complete proofs of Lemmas~\ref{Lemma3.2} and~\ref{Lemma3.3}, Theorem~\ref{Theorem3.19} and Corollary~\ref{Corollary3.20}.

Firstly, the complete proof of Lemma~\ref{Lemma3.2} is as follows.

\begin{proof}[Proof of Lemma~\ref{Lemma3.2}]
It is obvious that \smash{$\pi_{\operatorname{CP}^2 \times U}$} is a holomorphic map. Moreover, the rank of the Jacobian matrix of \smash{$\pi_{\operatorname{CP}^2 \times U}$} is equal to 1 at every point of \smash{$\operatorname{CP}^2\times U$}. In addition, $\forall t \in U$, \smash{$\pi_{\operatorname{CP}^2 \times U}^{-1}(t)=\operatorname{CP}^2\times \{ t \}$} is a compact complex submanifold of $\operatorname{CP}^2\times U$. So $\bigl(\operatorname{CP}^2\times U, U,\pi_{\operatorname{CP}^2\times U}\bigr)$ is a complex analytic family of $\operatorname{CP}^2$ over a 1-dimensional complex manifold $U$.

Now \smash{$\widetilde{\mathcal{S}}=\bigl\{ (x, \tau) \in \operatorname{CP}^2 \times U|x \in f_{\tau}(C_0(\tau)) \bigr\}$} is a complex submanifold of $\operatorname{CP}^2 \times U$ and the rank of the Jacobian matrix of \smash{$\pi_{\operatorname{CP}^2 \times U}|_{\widetilde{\mathcal{S}}}$} is equal to 1 at every point of \smash{$\widetilde{\mathcal{S}}$}. $\forall t \in U$, $\smash{\pi_{\operatorname{CP}^2 \times U}|_{\widetilde{\mathcal{S}}}^{-1}(t)}=f_{t}(C_0(t))\times \{ t \}$ is a compact complex submanifold of $\widetilde{\mathcal{S}}$. So \smash{$\bigl(\widetilde{\mathcal{S}},U,\pi_{\operatorname{CP}^2 \times U}|_{\widetilde{\mathcal{S}}}\bigr)$} is a complex analytic family of elliptic curves over the 1-dimensional complex manifold $U$.

Therefore, $\bigl(\widetilde{\mathcal{S}},U,\pi_{\operatorname{CP}^2 \times U}|_{\widetilde{\mathcal{S}}}\bigr)$ is a complex analytic family of elliptic curves and $\bigl(\operatorname{CP}^2\times U, U,$ $\pi_{\operatorname{CP}^2\times U}\bigr)$ is a complex analytic family of $\operatorname{CP}^2$. So Lemma~\ref{Lemma3.2} is proved. \end{proof}

Secondly, the proof of Lemma~\ref{Lemma3.3} was simply described in \cite{Takayuki-Takato1}. Here, a complete proof of Lemma~\ref{Lemma3.3} is as follows.

\begin{proof}[Proof of Lemma~\ref{Lemma3.3}]
For all $\widetilde{\tau} \in U$, let $N_{\widetilde{\tau}}$ be the normal bundle of the strict transform for the elliptic curve $f_{\widetilde{\tau}}(C_0(\widetilde{\tau}))\times \{ \widetilde{\tau} \} \cong C_0(\widetilde{\tau})$ after blowing up at nine points
\[\{ \widetilde{p}_1=(f_{\widetilde{\tau}}(p_1),\widetilde{\tau}),\dots , \widetilde{p}_9=(f_{\widetilde{\tau}}(p_9),\widetilde{\tau}) \}\subset f_{\widetilde{\tau}}(C_0(\widetilde{\tau}))\times \{ \widetilde{\tau} \}\]
with
\[ \{ p_1=f_{\widetilde{\tau},1}(\hat{p}_1), p_2=f_{\widetilde{\tau},2}(\hat{p}_2),\dots ,p_8=f_{\widetilde{\tau},8}(\hat{p}_8),p_9 \}\subset C_0(\widetilde{\tau})=\mathbb{C}/ \langle 1,\widetilde{\tau} \rangle\]
and $(\hat{p}_1,\dots ,\hat{p}_8)\in $$U_1 \times \dots \times U_8$. In addition, let $C(\widetilde{\tau})\times \{ \widetilde{\tau} \}$ be the strict transform of $ f_{\widetilde{\tau}}(C_0(\widetilde{\tau}))\times \{ \widetilde{\tau} \}$. Here, $N_{\widetilde{\tau}}$ is assumed to satisfy the Diophantine condition.

Then $N_{\widetilde{\tau}}$ is isomorphic to
\[
\mathcal{O}_{\operatorname{CP}^2}(3)|_{f_{\widetilde{\tau}}(C_0(\widetilde{\tau}))}\otimes \mathcal{O}_{f_{\widetilde{\tau}}(C_0(\widetilde{\tau}))}(-f_{\widetilde{\tau}}(p_1)-f_{\widetilde{\tau}}(p_2)-\dots -f_{\widetilde{\tau}}(p_9))\in \operatorname{Pic}^0(f_{\widetilde{\tau}}(C_0(\widetilde{\tau}))).
\]
Furthermore, taking use of notations from Section~\ref{sec2}, since $V_s$ is biholomorphic to a topologically trivial annulus bundle over the strict transform $C(\widetilde{\tau})$, $V_s$ is homotopic to $ S_1 \times S_2 \times S_3$ with~$ { \{S_1 , S_2 , S_3 \}}$ being circles and $ S_1 \times S_2$ being $C^\infty$ sections for $s\in \Delta \setminus \{ 0\}$ and sufficiently small \cite{Takayuki-Takato2} (also see Theorem~\ref{Theorem2.8}). Then through integrating the global holomorphic form $\sigma_s$ defined in Section~\ref{sec2} on $ S_1 \times S_2$, the result value is $q-p \cdot \widetilde{\tau}$ \cite{Takayuki-Takato1}. So there is an equation
\[9p_0-\sum _{j=1}^8p_j-p_9=q-p\cdot \widetilde{\tau} \quad \mod \langle 1, \widetilde{\tau} \rangle \]
with $p_0$ being an inflection point on $C_0(\widetilde{\tau})$ derived in \cite{Takayuki-Takato2}. Then for $\hat{t}=(\widetilde{\tau},\hat{p}_1,\dots ,\hat{p}_8)\in T$, $p_9=f_9\bigl(\hat{t}\bigr)\in C_0(\widetilde{\tau})=\mathbb{C}/ \langle 1, \widetilde{\tau} \rangle$.

Therefore, $N_{\widetilde{\tau}}$ satisfies the Diophantine condition corresponding to $(p,q)$ if and only if
\[9p_0-\sum _{j=1}^8p_j-p_9=q-p\cdot \widetilde{\tau} \quad \mod \langle 1, \widetilde{\tau} \rangle \]
with $p_0$ being an inflection point on $C_0(\widetilde{\tau})$.

In addition, taking use of the detailed definition of the map $\widetilde{\pi}$, $(\mathcal{F},T, \widetilde{\pi})$ and $(\mathcal{S},T, \widetilde{\pi}|_\mathcal{S})$ can be proved to be complex analytic families directly.

In conclusion, through blowing up nine points on every fiber of $\bigl(\operatorname{CP}^2 \times T,T, \pi\bigr)$ such that the normal bundle of the strict transform of the elliptic curve mentioned in Section~\ref{sec3.1} embedded in each fiber satisfies the Diophantine condition corresponding to the Diophantine number pair~$(p,q)$, the complex analytic families $(\mathcal{F},T, \widetilde{\pi})$ and $(\mathcal{S},T, \widetilde{\pi}|_\mathcal{S})$ can be constructed. So Lemma~\ref{Lemma3.3} is proved.
\end{proof}

Thirdly, the complete proof of Theorem~\ref{Theorem3.19} is as follows.

\begin{proof}[Proof of Theorem~\ref{Theorem3.19}]
First of all, it is to prove that there exists $b>0$ such that the metric $h \cdot {\rm e}^{-b \Psi_s}$ is a~smooth hermitian metric on $L|_{S(\tau)\backslash C(\tau)}$ with $\Theta_{h \cdot {\rm e}^{-b \Psi_s}}>0$.

On the set $S(\tau)\backslash W$, since $\Psi_s$ is constant, $\partial \bar{\partial}\Psi_s=0$. Then
\[\Theta_{h \cdot {\rm e}^{-b \Psi_s}}=\Theta_h + \sqrt{-1} \frac{b}{\pi} \partial \bar{\partial} \Psi_s =\Theta_h=\Theta_{h_L}>0.\]

On the set $\{ [(z,w)]\in W\backslash C(\tau) \mid r_2 \leq |w|<r \}$, from the construction of $\Psi_s$ above, since~$\widetilde{\theta}_s(x)$ is smooth for $x>0$, the coefficients of $\partial \bar{\partial}\Psi_s$ are bounded. Moreover, since $\Theta_{h_L}>0$ on~$S(\tau)\backslash C(\tau)$, the coefficients of $\Theta_h$ have a positive lower bound on $\{ [(z,w)]\in W\backslash C(\tau) \mid r_2 \leq |w|<r \}$. Then $\exists$ sufficiently small $b>0$ such that $\Theta_{h \cdot {\rm e}^{-b \Psi_s}}=\Theta_h + \sqrt{-1} \frac{b}{\pi}\partial \bar{\partial} \Psi_s >0$.

On the set $\{ [(z,w)]\in W\backslash C(\tau) \mid |w|<r_2 \}$,
$\sqrt{-1} \partial \bar{\partial} \Psi_s=2 \sqrt{-1} \cdot {\rm d}w \wedge {\rm d}\bar{w} /|w|^2>0$.
So $\Theta_{h \cdot {\rm e}^{-b \Psi_s}}=\Theta_h + \sqrt{-1} \frac{b}{\pi}\partial \bar{\partial} \Psi_s >0$.

The next step is to prove that $\Theta_{h \cdot {\rm e}^{-b \Psi_s}}$ is of the form
\[
\omega|_{W_{\varepsilon_0}}=\frac{ (L.C(\tau))}{\operatorname{ Im}\tau}\cdot \sqrt{-1} {\rm d}z \wedge {\rm d} \bar{z}+\frac{2b}{\pi} \cdot \frac{\sqrt{-1}{\rm d}w \wedge {\rm d} \bar{w}}{|w|^2}
\]
on the set $W_{\varepsilon_0}=\{[(z,w)]\in W \backslash C(\tau)\mid |w|<\sqrt{\varepsilon_0}r\}$.

On the set $W_{\varepsilon_0}=\{[(z,w)]\in W \backslash C(\tau)\mid |w|<\sqrt{\varepsilon_0}r\}$,
\begin{align*}
\Theta_{h \cdot {\rm e}^{-b \Psi_s}}&=\Theta_{(\pi_W)^*h_{C(\tau)}}+\frac{2b}{\pi}\cdot \frac{\sqrt{-1}{\rm d}w \wedge {\rm d} \bar{w}}{|w|^2}\\
&=\frac{(L.C(\tau))}{\operatorname{ Im}\tau}\cdot \sqrt{-1} {\rm d}z \wedge {\rm d} \bar{z}+\frac{2b}{\pi}\cdot \frac{\sqrt{-1}{\rm d}w \wedge {\rm d} \bar{w}}{|w|^2}.
\end{align*}
 So the curvature form $\Theta_{h \cdot {\rm e}^{-b\Psi_s}}$ gives a complete K\"ahler metric on $S(\tau)\backslash C(\tau)$.

The last step is to prove that $\omega|_{W_{\varepsilon_0}}$ is Ricci-flat on the set $W_{\varepsilon_0}=\{ [(z,w)]\in W \backslash C(\tau)\mid |w|<\sqrt{\varepsilon_0}r \}$.

The coefficient matrix of $\omega|_{W_{\varepsilon_0}}$ on the set $W_{\varepsilon_0}=\{ [(z,w)]\in W \backslash C(\tau)\mid |w|<\sqrt{\varepsilon_0}r \}$ is \[(g_{ij})_{2 \times 2}= \begin{pmatrix}
 \frac{2 (L.C(\tau))}{\operatorname{ Im}\tau} &0\\
 0 & \frac{4b}{\pi |w|^2}
\end{pmatrix}.\]
Then the determinate of the coefficient matrix is
\[\det((g_{ij})_{2\times 2})=\frac{8b (L.C(\tau))}{\pi \operatorname{ Im}\tau \cdot |w|^2}.\]
So $\operatorname{Ric}(\omega|_{W_{\varepsilon_0}})=0$. Therefore, $\omega|_{W_{\varepsilon_0}}$ is Ricci-flat.

In conclusion, $\exists \ b>0$ such that $\Theta_{h\cdot {\rm e}^{-b\Psi_s}}$ gives a complete K\"ahler metric on $S(\tau)\backslash C(\tau)$. In addition, $\Theta_{h\cdot {\rm e}^{-b\Psi_s}}$ is of the form
\[
\omega|_{W_{\varepsilon_0}}=\frac{(L.C(\tau))}{\operatorname{ Im}\tau}\cdot \sqrt{-1} {\rm d}z \wedge {\rm d} \bar{z}+\frac{2b}{\pi} \cdot \frac{\sqrt{-1}{\rm d}w \wedge {\rm d} \bar{w}}{|w|^2}
\]
on the set $W_{\varepsilon_0}=\{[(z,w)]\in W \backslash C(\tau)\mid |w|<\sqrt{\varepsilon_0}r\}$. Furthermore, $\omega|_{W_{\varepsilon_0}}$is Ricci-flat. That is to say, $\omega|_{W_{\varepsilon_0}}$ corresponds to a Ricci-flat K\"ahler metric. So Theorem~\ref{Theorem3.19} is proved.
\end{proof}

Finally, the complete proof of Corollary~\ref{Corollary3.20} is as follows.

\begin{proof}[Proof of Corollary~\ref{Corollary3.20}]
Since $\widetilde{\pi}|_{\mathcal{F}\backslash \mathcal{S}}$ is a holomorphic surjective submersion and $f$ is a biholomorphic map from~$D_3$ to $f(D_3)$, then $f \circ \widetilde{\pi}|_{\mathcal{F}\backslash \mathcal{S}}$ is also a holomorphic surjective submersion. $(f \circ \widetilde{\pi}|_{\mathcal{F}\backslash \mathcal{S}})^{-1}(0)=j(S\backslash C)$.

Moreover, since $\widetilde{\pi}|_{\mathcal{F}\backslash \mathcal{S}}$ is locally trivial in the $C^\infty$ category and $f$ is a diffeomorphism, then \smash{$f \circ \widetilde{\pi}|_{\mathcal{F}\backslash \mathcal{S}}$} is also locally trivial in the $C^\infty$ category.

In conclusion, $f \circ \widetilde{\pi}|_{\mathcal{F}\backslash \mathcal{S}}\colon \mathcal{F} \rightarrow f(D_3)$ is a deformation family of the non-compact complete K\"ahler manifold $j(S\backslash C)$ over a 9-dimensional complex manifold $f(D_3)$ satisfying the conditions in Definition~\ref{Definition3.5}. So Corollary~\ref{Corollary3.20} is proved.
\end{proof}


\begin{thebibliography}{99}
\footnotesize\itemsep=0pt

\bibitem{Anna}
Abasheva A., D\'eev R., Complex surfaces with many algebraic structures,
 \href{https://doi.org/10.1093/imrn/rnad190}{\textit{Int. Math. Res. Not.}}
 \textbf{2024} (2024), 7379--7400,
 \href{http://arxiv.org/abs/2303.10764}{arXiv:2303.10764}.

\bibitem{Yukitaka-Klaus}
Abe Y., Kopfermann K., Toroidal groups. Line bundles, cohomology and
 quasi-abelian varieties, \textit{Lecture Notes in Math.}, Vol.~1759,
 \href{https://doi.org/10.1007/b80605}{Springer}, Berlin, 2001.

\bibitem{Vladimir}
Arnold V.I., Bifurcations of invariant manifolds of differential equations, and
 normal forms of neighborhoods of elliptic curves,
 \href{https://doi.org/10.1007/BF01076024}{\textit{Funct. Anal. Appl.}}
 \textbf{10} (1976), 249--259.

\bibitem{Edoardo-Elizabeth-Francisco}
Ballico E., Gasparim E., Rubilar F., 20~open questions about deformations of
 compactifiable manifolds,
 \href{https://doi.org/10.1007/s40863-021-00213-8}{\textit{S\~{a}o Paulo~J.
 Math. Sci.}} \textbf{15} (2021), 661--681,
 \href{http://arxiv.org/abs/2004.11299}{arXiv:2004.11299}.

\bibitem{Marco}
Brunella M., On {K}\"ahler surfaces with semipositive {R}icci curvature,
 \textit{Riv. Math. Univ. Parma~(N.S.)} \textbf{1} (2010), 441--450.

\bibitem{Jean}
Demailly J.-P., Complex analytic and differential geometry, Universit\`e de
 Grenoble~I, Grenoble, 1997.

\bibitem{Lawrence}
Evans L.C., Partial differential equations, \textit{Grad. Stud. Math.},
 Vol.~19, \href{https://doi.org/10.1090/gsm/019}{American Mathematical
 Society}, Providence, RI, 1998.

\bibitem{Yoshio}
Fujimoto Y., On rational elliptic surfaces with multiple fibers,
 \href{https://doi.org/10.2977/prims/1195171661}{\textit{Publ. Res. Inst.
 Math. Sci.}} \textbf{26} (1990), 1--13.

\bibitem{Elizabeth-Francisco}
Gasparim E., Rubilar F., Deformations of noncompact {C}alabi-{Y}au manifolds,
 families and diamonds, in Geometry at the {F}rontier---{S}ymmetries and
 {M}oduli {S}paces of {A}lgebraic {V}arieties, \textit{Contemp. Math.}, Vol.~766, \href{https://doi.org/10.1090/conm/766/15378}{American Mathematical
 Society}, Providence, RI, 2021, 117--132,
 \href{http://arxiv.org/abs/1908.09045}{arXiv:1908.09045}.

\bibitem{Phillip}
Griffiths P., Harris J., Principles of algebraic geometry, \textit{Wiley Classics Lib.},
 \href{https://doi.org/10.1002/9781118032527}{John Wiley \& Sons}, New York,
 1994.

\bibitem{Vincent-Ahmed}
Guedj V., Zeriahi A., Degenerate complex {M}onge--{A}mp\`ere equations,
 \textit{EMS Tracts Math.}, Vol.~26,
 \href{https://doi.org/10.4171/167}{European Mathematical Society (EMS)},
 Z\"urich, 2017.

\bibitem{Robin}
Hartshorne R., Algebraic geometry, \textit{Grad. Texts in Math.}, Vol.~52,
 \href{https://doi.org/10.1007/978-1-4757-3849-0}{Springer}, New York, 1977.

\bibitem{Hans}
Hein H.J., Complete Calabi--Yau metrics
 from~$\mathbb{P}^2\#9\overline{\mathbb{P}}^2$,
 \href{http://arxiv.org/abs/1003.2646}{arXiv:1003.2646}.

\bibitem{Kunihiko}
Kodaira K., Complex manifolds and deformation of complex structures, \textit{Class.
 Math.}, \href{https://doi.org/10.1007/b138372}{Springer}, Berlin, 2005.

\bibitem{Takayuki-Takato1}
Koike T., Uehara T., A gluing construction of projective~{K}3 surfaces,
 \href{http://arxiv.org/abs/1903.01444}{arXiv:1903.01444}.

\bibitem{Takayuki-Takato2}
Koike T., Uehara T., A gluing construction of {K}3 surfaces,
 \href{https://doi.org/10.46298/epiga.2022.volume6.8504}{\textit{\'Epijournal
 G\'eom. Alg\'ebrique}} \textbf{6} (2022), 12, 15~pages,
 \href{http://arxiv.org/abs/2108.07168}{arXiv:2108.07168}.

\bibitem{Julius-David}
Ross J., Nystr\"om D.W., Homogeneous {M}onge--{A}mp\`ere equations and
 canonical tubular neighbourhoods in {K}\"ahler geometry,
 \href{https://doi.org/10.1093/imrn/rnw200}{\textit{Int. Math. Res. Not.}}
 \textbf{2017} (2017), 7069--7108,
 \href{http://arxiv.org/abs/1403.3282}{arXiv:1403.3282}.

\bibitem{Alexandre}
Sukhov A., Regularized maximum of strictly plurisubharmonic functions on an
 almost complex manifold,
 \href{https://doi.org/10.1142/S0129167X13500973}{\textit{Internat.~J. Math.}}
 \textbf{24} (2013), 1350097, 6~pages,
 \href{http://arxiv.org/abs/1303.5312}{arXiv:1303.5312}.

\bibitem{Tomasz-Halszka}
Szemberg T., Tutaj-Gasi\'nska H., General blow-ups of the projective plane,
 \href{https://doi.org/10.1090/S0002-9939-02-06488-2}{\textit{Proc. Amer.
 Math. Soc.}} \textbf{130} (2002), 2515--2524.

\bibitem{Tetsuo}
Ueda T., On the neighborhood of a compact complex curve with topologically
 trivial normal bundle,
 \href{https://doi.org/10.1215/kjm/1250521670}{\textit{J.~Math. Kyoto Univ.}}
 \textbf{22} (1982), 583--607.

\bibitem{Christian}
Vogt C., Line bundles on toroidal groups,
 \href{https://doi.org/10.1515/crll.1982.335.197}{\textit{J.~Reine Angew.
 Math.}} \textbf{335} (1982), 197--215.

\bibitem{Claire}
Voisin C., Hodge theory and complex algebraic geometry.~{I}, \textit{Cambridge
 Stud. Adv. Math.}, Vol.~76,
 \href{https://doi.org/10.1017/CBO9780511615344}{Cambridge University Pres}s,
 Cambridge, 2002.

\end{thebibliography}
\end{document}
